\chapter{多元微分与场论}\label{ch:c3}

温度沿轨迹的变化率、平衡位置附近的势能，以及电场的局部结构，都依赖多个坐标。偏导数描述固定其他坐标时的变化，全微分则给出任意小位移下的线性变化。存在约束时，允许的一阶位移还须留在切空间内。

本讲依次讨论偏导与可微、无约束和约束极值，以及梯度、散度、旋度的几何来源。球坐标极角从正$z$轴量起，方位角从正$x$轴朝正$y$轴增加。向量分量均按当地单位基矢定义。势与场的关系还需结合定义域与边界条件使用。


\section{偏导数、全微分与复合函数}\label{sec:c3-diff}
多个状态量同时变化时，固定变量的选择会改变偏导结果。全微分与链式法则沿明确的依赖关系计算，物理约束也随之进入求导。

\subsection{偏导数与独立变量}
对$f(x,y,z)$求$x$的偏导数，是在$y,z$保持不变时取一元差商的极限：
\[f_x(x,y,z)=\lim_{h\to0}\frac{f(x+h,y,z)-f(x,y,z)}h.\]
偏导数必须注明固定的变量。例如气体的$(\partial p/\partial V)_T$与$(\partial p/\partial V)_S$对应不同过程，一般具有不同数值。独立变量选定后，求导时便按这一约定保持其他变量不变。

若函数在一点附近可微，则增量能写成
\[\Delta f=f_x\Delta x+f_y\Delta y+f_z\Delta z+o\!\left(\sqrt{\Delta x^2+\Delta y^2+\Delta z^2}\right).
\]
其中线性部分称为全微分：$\dd f=f_x\dd x+f_y\dd y+f_z\dd z$。各偏导存在并不自动保证可微；一个常用的充分条件是各一阶偏导在该点附近连续。偏导是沿坐标轴的变化率，全微分则用同一个线性映射同时控制所有方向的微小变化。

\subsection{从坐标方向到任意小位移}
偏导数只沿坐标轴取极限，因此知道$f_x,f_y$仅掌握了两条特殊路径。可微性要求同一个线性式$f_x\Delta x+f_y\Delta y$对任意趋近方向都成立，而且余项除以位移长度后趋零。这是方向无关的控制条件，比逐条轴线上的差商极限更强。

一阶偏导连续为何足以保证可微，可以把二维增量拆成两步：
\[f(x+h,y+k)-f(x,y)
=\{f(x+h,y+k)-f(x,y+k)\}+\{f(x,y+k)-f(x,y)\}.\]
对两项分别用一元中值定理，得到$f_x(\xi,y+k)h+f_y(x,\eta)k$。偏导连续使系数趋向$f_x(x,y),f_y(x,y)$，剩余误差为$o(\sqrt{h^2+k^2})$。习题~\ref{ex:c3-h3}则说明，如果只有轴向偏导存在，这一步的系数连续性可能失效。

可微时，一阶近似的几何对象是切平面。对$z=f(x,y)$，在$(x_0,y_0)$附近有$z-z_0=f_x\Delta x+f_y\Delta y$；梯度的两个分量给出沿坐标方向的坡度，其他方向的坡度由同一个平面确定。若某函数沿不同路径呈现无法由同一平面描述的一阶变化，就不可微。

\subsection{复合函数的链式法则}
若$w=f(x,y)$，而$x=x(s,t),y=y(s,t)$，则固定$t$沿$s$变化时
\[w_s=f_xx_s+f_yy_s,\qquad w_t=f_xx_t+f_yy_t.\]
右侧的$f_x,f_y$应在当前$(x(s,t),y(s,t))$取值。若再有显含时间的场$f(x,y,z,t)$，沿轨迹的变化率为
\[\frac{\dd f}{\dd t}=\frac{\partial f}{\partial t}+\vct v\cdot\nabla f.\]
第一项描述固定位置处场随时间变化，第二项描述运动者穿过空间不均匀场而引起的变化。

链式法则可直接从这一近似推出。令$x=x(s,t),y=y(s,t)$，固定$t$改变$s$时，$(\Delta x,\Delta y)=(x_s,y_s)\Delta s+o(\Delta s)$，代入全微分即得$w_s=f_xx_s+f_yy_s$。参数曲线的速度没有单位长度归一化，所以这个参数导数与单位方向导数还相差一个速率因子。

\subsection{隐函数的分母与局部可解性}
设$F(x,y,z)=0$在一点满足$F_z\ne0$，并具备连续偏导，则局部可将$z$看作$x,y$的函数。沿约束面有
\[F_x\dd x+F_y\dd y+F_z\dd z=0,\qquad
z_x=-F_x/F_z,\quad z_y=-F_y/F_z.
\]
当$F_z=0$，这些商式失效；曲面本身可能仍然光滑，只是不能在该处以$z$作为因变量。

隐函数方程$F(x,y,z)=0$把三个坐标限制在同一曲面上。若$F_z\ne0$，给定足够小的$\Delta x,\Delta y$后，改变$z$能在一阶上抵消$F$的变化，所需增量为$\Delta z\simeq-(F_x\Delta x+F_y\Delta y)/F_z$。分母非零使这种调节局部唯一。

分母为零时不一定是曲面本身奇异。如果$F_x\ne0$，仍可改写$x=x(y,z)$；此时曲面只是对$xy$平面的投影发生退化。只有$\nabla F=0$时，通常的正规水平面论证才整体失效，还须用原方程或参数化判断几何。例~\ref{ex:c3-e2}中应先计算三个独立偏导，再按所选因变量取商，避免把$z$尚未求出的依赖关系提前混入$F_x$。

混合偏导的次序也要区分。连续二阶偏导保证$f_{xy}=f_{yx}$，并非任意二次偏导在存在时都自动相等。光滑势函数、解析场等竞赛常见模型满足这一条件，所以可以交换；在分段定义或奇点附近，仍应回到定义验证。

\begin{example}[偏导数与混合偏导]\label{ex:c3-e1}
求$f=x^2+y^2/4+\sin z$的一阶偏导；求$g=\ee^{y\sin x}$的$g_x,g_y$；对$h=\ee^{x^2y}$分别计算$h_{xy},h_{yx}$。 见图~\ref{fig:c3-e1-q}。
\end{example}

\begin{figure}[H]\centering
\begin{tikzpicture}[>=Stealth,line cap=round,line join=round,x=2cm,y=1.1cm]
\draw[axis](-.2,0)--(1.9,0)node[right]{$x$};
\draw[axis](0,-.2)--(0,3.2)node[above]{$f$};
\draw[curve,domain=-.15:1.65,samples=101]plot(\x,{(\x)^2});
\draw[dashed,domain=.55:1.55]plot(\x,{2*\x-1});
\draw[guide](1,0)--(1,1)--(0,1);
\fill(1,1)circle(1.5pt)node[above left]{$P=(1,1)$};
\node[below]at(1,0){$1$};\node[left]at(0,1){$1$};
\node[right,align=left]at(2,.95){$y_0=z_0=0$\\[5pt]$f=x^2$\\[5pt]$f_x(P)=2$};
\end{tikzpicture}
\caption{取题设函数在$y_0=z_0=0$的真实截线；虚线为$P$点切线$f=2x-1$。}\label{fig:c3-e1-q}
\end{figure}
\begin{solution}
偏导时先固定其余独立坐标。图~\ref{fig:c3-e1-q}固定$y=z=0$，截线为$f=x^2$，在$P$点的斜率为$2$。若其他坐标也随$x$变化，轨迹导数还含相应链式项。三组函数分别检验多项式分项求导、复合指数链式法则，以及乘积中的两个因子是否都被求导。


按图~\ref{fig:c3-e1-s}中所示的运算路径，求$x$导数时$y,z$作为常数，故
\[f_x=2x,\qquad f_y=y/2,\qquad f_z=\cos z.\]
复合指数函数需要先对外层求导，再对指数求导：
\[g_x=y\cos x\,\ee^{y\sin x},\qquad
 g_y=\sin x\,\ee^{y\sin x}.\]
对$h$先求$h_x=2xy\ee^{x^2y}$，再对$y$求导时，两因子都依赖$y$：
\[h_{xy}=2x\ee^{x^2y}+2x^3y\ee^{x^2y}
=\boxed{2x(1+x^2y)\ee^{x^2y}}.\]
反序求导从$h_y=x^2\ee^{x^2y}$出发：
\[h_y=x^2\ee^{x^2y},\qquad
h_{yx}=2x\ee^{x^2y}+x^2\ee^{x^2y}(2xy).\]
与先对$x$求导后的结果逐项一致。若只对指数求导，会漏掉第一项$2x\ee^{x^2y}$；若把$y$误当作$x$的函数，又会引入题目未给出的$y'$。

图~\ref{fig:c3-e1-s}用两条求导路径汇到同一个混合导数。两条路径的中间函数不同，最终相等依赖本题的光滑性。分别写出两个中间函数，便可核对混合求导中是否遗漏乘积项。

\begin{figure}[H]\centering
\begin{tikzpicture}[>=Stealth,line cap=round,line join=round]
\node at(0,0){$h$};\node at(3,1){$h_x=2xy\ee^{x^2y}$};\node at(3,-1){$h_y=x^2\ee^{x^2y}$};\node at(7,0){$h_{xy}=h_{yx}$};\draw[axis](.5,.2)--(1.5,.8)node[midway,above]{$\partial_x$};\draw[axis](.5,-.2)--(1.5,-.8)node[midway,below]{$\partial_y$};\draw[axis](4.8,.8)--(6,.2)node[midway,above]{$\partial_y$};\draw[axis](4.8,-.8)--(6,-.2)node[midway,below]{$\partial_x$};
\end{tikzpicture}
\caption{两条求导路径各自使用乘积与链式法则，再比较结果。}\label{fig:c3-e1-s}
\end{figure}
\end{solution}

\Needspace{8\baselineskip}
\begin{example}[隐式曲面的局部斜率]\label{ex:c3-e2}
曲面$3x^2z-x^2y^2+2z^3+3yz-5=0$上，在能将$z$表示为$x,y$函数的区域，求$z_x,z_y$，并给出适用条件。 见图~\ref{fig:c3-e2-q}。
\end{example}

\begin{figure}[H]\centering
\begin{tikzpicture}[>=Stealth,line cap=round,line join=round]
\draw[axis](-.4,0)--(3,0)node[right]{$y$};\draw[axis](0,-.2)--(0,2)node[above]{$z$};\draw[curve]plot coordinates{(2.5377778,0.6000000) (2.5090183,0.6053333) (2.4806485,0.6106667) (2.4526570,0.6160000) (2.4250334,0.6213333) (2.3977671,0.6266667) (2.3708481,0.6320000) (2.3442669,0.6373333) (2.3180140,0.6426667) (2.2920805,0.6480000) (2.2664574,0.6533333) (2.2411365,0.6586667) (2.2161095,0.6640000) (2.1913684,0.6693333) (2.1669057,0.6746667) (2.1427137,0.6800000) (2.1187854,0.6853333) (2.0951138,0.6906667) (2.0716920,0.6960000) (2.0485136,0.7013333) (2.0255720,0.7066667) (2.0028613,0.7120000) (1.9803753,0.7173333) (1.9581083,0.7226667) (1.9360546,0.7280000) (1.9142088,0.7333333) (1.8925654,0.7386667) (1.8711194,0.7440000) (1.8498657,0.7493333) (1.8287994,0.7546667) (1.8079158,0.7600000) (1.7872103,0.7653333) (1.7666784,0.7706667) (1.7463157,0.7760000) (1.7261180,0.7813333) (1.7060811,0.7866667) (1.6862011,0.7920000) (1.6664740,0.7973333) (1.6468961,0.8026667) (1.6274636,0.8080000) (1.6081729,0.8133333) (1.5890205,0.8186667) (1.5700031,0.8240000) (1.5511171,0.8293333) (1.5323595,0.8346667) (1.5137270,0.8400000) (1.4952165,0.8453333) (1.4768251,0.8506667) (1.4585498,0.8560000) (1.4403878,0.8613333) (1.4223362,0.8666667) (1.4043923,0.8720000) (1.3865535,0.8773333) (1.3688172,0.8826667) (1.3511809,0.8880000) (1.3336420,0.8933333) (1.3161982,0.8986667) (1.2988472,0.9040000) (1.2815865,0.9093333) (1.2644140,0.9146667) (1.2473275,0.9200000) (1.2303249,0.9253333) (1.2134040,0.9306667) (1.1965628,0.9360000) (1.1797993,0.9413333) (1.1631115,0.9466667) (1.1464976,0.9520000) (1.1299557,0.9573333) (1.1134839,0.9626667) (1.0970804,0.9680000) (1.0807436,0.9733333) (1.0644716,0.9786667) (1.0482629,0.9840000) (1.0321158,0.9893333) (1.0160287,0.9946667) (1.0000000,1.0000000) (0.9840282,1.0053333) (0.9681118,1.0106667) (0.9522493,1.0160000) (0.9364393,1.0213333) (0.9206803,1.0266667) (0.9049711,1.0320000) (0.8893102,1.0373333) (0.8736963,1.0426667) (0.8581281,1.0480000) (0.8426044,1.0533333) (0.8271239,1.0586667) (0.8116854,1.0640000) (0.7962876,1.0693333) (0.7809295,1.0746667) (0.7656099,1.0800000) (0.7503276,1.0853333) (0.7350815,1.0906667) (0.7198706,1.0960000) (0.7046938,1.1013333) (0.6895500,1.1066667) (0.6744383,1.1120000) (0.6593576,1.1173333) (0.6443069,1.1226667) (0.6292854,1.1280000) (0.6142919,1.1333333) (0.5993257,1.1386667) (0.5843858,1.1440000) (0.5694713,1.1493333) (0.5545813,1.1546667) (0.5397149,1.1600000) (0.5248714,1.1653333) (0.5100499,1.1706667) (0.4952496,1.1760000) (0.4804696,1.1813333) (0.4657092,1.1866667) (0.4509676,1.1920000) (0.4362441,1.1973333) (0.4215379,1.2026667) (0.4068483,1.2080000) (0.3921745,1.2133333) (0.3775159,1.2186667) (0.3628718,1.2240000) (0.3482414,1.2293333) (0.3336242,1.2346667) (0.3190194,1.2400000) (0.3044264,1.2453333) (0.2898445,1.2506667) (0.2752732,1.2560000) (0.2607119,1.2613333) (0.2461598,1.2666667) (0.2316165,1.2720000) (0.2170814,1.2773333) (0.2025538,1.2826667) (0.1880332,1.2880000) (0.1735191,1.2933333) (0.1590108,1.2986667) (0.1445079,1.3040000) (0.1300099,1.3093333) (0.1155162,1.3146667) (0.1010263,1.3200000) (0.0865396,1.3253333) (0.0720558,1.3306667) (0.0575743,1.3360000) (0.0430947,1.3413333) (0.0286164,1.3466667) (0.0141390,1.3520000) (-0.0003380,1.3573333) (-0.0148150,1.3626667) (-0.0292924,1.3680000) (-0.0437707,1.3733333) (-0.0582504,1.3786667) (-0.0727317,1.3840000) (-0.0872153,1.3893333) (-0.1017014,1.3946667) (-0.1161905,1.4000000)};\fill(1,1)circle(1.5pt)node[above right]{$P(0,1,1)$};\node at(2.6,1.8){$x=0$};
\end{tikzpicture}
\caption{原曲面在$x=0$的截线经过$P$，此处$z_y=-1/3$。}\label{fig:c3-e2-q}
\end{figure}
\begin{solution}
曲面方程中的$x,y,z$在计算$F$的偏导时先视为独立变量，随后才用约束令$z$依赖$x,y$。图~\ref{fig:c3-e2-q}取$x=0$截面上的具体点，展示“局部斜率”如何从隐式方程得到，而非预先解出整个三次曲面的所有分支。


记左侧为$F$。先在$x,y,z$彼此独立时计算
\[F_x=6xz-2xy^2,\quad F_y=-2x^2y+3z,\quad F_z=3x^2+6z^2+3y.\]
沿曲面固定 $y$，图~\ref{fig:c3-e2-s}中的增量满足 $F_x\dd x+F_z\dd z=0$；除以 $\dd x$ 即 $F_x+F_zz_x=0$。固定 $x$ 同理有 $F_y+F_zz_y=0$，因此
\[\boxed{z_x=-\frac{6xz-2xy^2}{3x^2+6z^2+3y},\qquad
z_y=-\frac{-2x^2y+3z}{3x^2+6z^2+3y}}.\]
这些式子要求分母非零。若分母为零而$F_x$或$F_y$非零，可以尝试改用$x$或$y$作为因变量；若三个偏导同时为零，则还须考察曲面的奇点结构。

在$P=(0,1,1)$，$F=0$且$\nabla F=(0,3,9)$，所以$z_x=0,z_y=-1/3$。切平面的法向为$\nabla F$，方程为$3\Delta y+9\Delta z=0$，即$z-1=-(y-1)/3$。图~\ref{fig:c3-e2-s}中沿固定$x$的截线有$\dd z=-\dd y/3$，而沿$x$方向的一阶变化为零。


例如 $3x^2z(x,y)$ 对 $x$ 求导时，两项分别为 $6xz$ 和 $3x^2z_x$；$2z^3$ 的导数为 $6z^2z_x$。这些包含 $z_x$ 的项收集后，系数正是上面独立求导得到的 $F_z$。

切平面只描述一阶变化，有限位移还受曲面弯曲影响。

\begin{figure}[H]\centering
\begin{tikzpicture}[>=Stealth,line cap=round,line join=round]
\draw[axis](-.2,0)--(3.6,0)node[right]{$y$};\draw[axis](0,-.2)--(0,2.6)node[above]{$z$};\draw[curve](.3,2.3)--(3.3,1.3);\fill(1.8,1.8)circle(1.5pt)node[above]{$P$};\draw[guide](1.8,1.8)--(2.7,1.8)--(2.7,1.5);\node[above]at(2.25,1.8){$\dd y$};\node[right]at(2.7,1.65){$\dd z$};\node at(1.8,.5){$\dd z=-\dd y/3$};
\end{tikzpicture}
\caption{在例示点沿固定$x$的截线，切线斜率为$-1/3$。}\label{fig:c3-e2-s}
\end{figure}
\end{solution}

\section{无约束极值与约束极值}\label{sec:c3-extrema}
受约束系统只能沿允许方向改变构形。求驻点后，还须分别检查可行边界及势能在可行方向上的变化，不能用无约束的二阶判据代替。

\subsection{驻点不等于极值点}
可微函数在定义域内部取局部极值时必须满足$\nabla f=0$。对二元函数，在驻点处记$A=f_{xx},B=f_{xy},C=f_{yy}$，二阶增量为
\[\Delta f=\tfrac12(A\Delta x^2+2B\Delta x\Delta y+C\Delta y^2)+o(\Delta x^2+\Delta y^2).\]
若$AC-B^2>0$且$A>0$，二次型正定，是严格极小值；若同一行列式为正且$A<0$，是严格极大值；若行列式为负，是鞍点。等于零时二阶项不足以判定，须考察更高阶或沿不同路径比较。定义域的边界点也须单独处理。

\subsection{二次型与Hessian判据}
在二元驻点附近，二次主部为$Q=Ah^2+2Bhk+Ck^2$。若$A\ne0$，配方为
\[Q=A(h+Bk/A)^2+\frac{AC-B^2}{A}k^2.\]
当$AC-B^2>0$时，两项系数与$A$同号，因此$A>0$对应严格极小，$A<0$对应严格极大。行列式为负时，两个主方向的曲率异号，总能找到使函数升高和降低的位移，故为鞍点。若$A=0$而$B\ne0$，行列式直接为负，同样是鞍点。

行列式为零时，至少有一个方向的二次项消失。沿这个方向，三阶、四阶乃至更高阶项可能决定符号，所以二阶判据只给出“不足以判断”，不直接等于“没有极值”。习题~\ref{ex:c3-h4}把三种四次函数并列，就是为了让相同的零Hessian与不同的真实极值性质分开。

三维及更高维中，Hessian为实对称矩阵，可用正交变换化为主曲率之和$Q=\sum_i\lambda_i\xi_i^2$。所有特征值为正才给出严格二次势阱；正负并存则为鞍点；有零特征值时还需分析零模态。对于机械势能，质量矩阵也影响小振动频率，但平衡处势能是否沿某方向下降，可先由这个二次型判断。

图~\ref{fig:c3-hessian-geometry}表示非退化驻点附近的二次主部。正定二次型的等值线围住驻点；不定二次型在不同方向上异号，同一点附近同时有上升与下降。退化时还需保留高阶项，不能仅由等值线示意图作结论。
\begin{figure}[H]\centering
\begin{tikzpicture}[>=Stealth,line cap=round,line join=round]
\begin{scope}
\draw[axis](-1.8,0)--(1.9,0)node[right]{$h_1$};\draw[axis](0,-1.4)--(0,1.5)node[above]{$h_2$};
\draw[curve](0,0)ellipse(.55 and .36);\draw[curve](0,0)ellipse(1 and .66);\draw[curve](0,0)ellipse(1.5 and 1);\fill(0,0)circle(1.5pt);\node at(0,-1.9){正定：各方向均增加};
\end{scope}
\begin{scope}[shift={(5.5,0)}]
\draw[axis](-1.8,0)--(1.9,0)node[right]{$h_1$};\draw[axis](0,-1.4)--(0,1.5)node[above]{$h_2$};
\draw[guide](-1.3,-1.3)--(1.3,1.3) (-1.3,1.3)--(1.3,-1.3);
\draw[curve,domain=-.95:.95,samples=55]plot({sqrt(.4+(\x)^2)},\x);\draw[curve,domain=-.95:.95,samples=55]plot({-sqrt(.4+(\x)^2)},\x);
\draw[curve,domain=-1.05:1.05,samples=55]plot(\x,{sqrt(.25+(\x)^2)});\draw[curve,domain=-1.05:1.05,samples=55]plot(\x,{-sqrt(.25+(\x)^2)});
\node at(1.45,.35){$+$};\node at(.25,1.15){$-$};\fill(0,0)circle(1.5pt);\node at(0,-1.9){不定：存在正负两类方向};
\end{scope}
\end{tikzpicture}
\caption{驻点附近的二次型几何；右图虚线表示二次主部为零的方向。}\label{fig:c3-hessian-geometry}
\end{figure}

\subsection{乘子法的几何含义}
在光滑约束$g(\vct r)=0$上，可行微位移满足$\nabla g\cdot\dd\vct r=0$。极值处若$\nabla f$沿某个切向有非零分量，便可沿该方向继续增减$f$；故正规约束下必有$\nabla f=\lambda\nabla g$。多个独立约束时，梯度属于各约束法向张成的空间：
\[\nabla f=\sum_j\lambda_j\nabla g_j,\qquad g_j=0.\]
可令$\mathcal L=f-\sum_j\lambda_jg_j$，对原变量与各乘子一并求偏导。对乘子的方程正好恢复原约束；只对原变量求导会得到欠定系统。

乘子法给出候选点。最后仍需比较函数值，检查边界、约束梯度不独立的点，以及消元中曾除去的因子为零的分支。若约束集合紧且函数连续，全局最大最小值存在，这时完整列出候选点后比较才足以得出全局结论。

\Needspace{8\baselineskip}
\subsection{正规约束、边界与候选点的完整性}
一条正规等值曲线$g=0$在某点满足$\nabla g\ne0$。允许的一阶位移垂直$\nabla g$；若目标梯度还有切向分量，便可沿曲线向一个方向使目标增加，向反方向使其减小，所以极值点的目标梯度只能沿约束法向。这给出乘子方程的几何必然性。

有两个约束时，$\nabla g_1,\nabla g_2$应线性独立，它们的叉积给出交线切向。极值处$\nabla f$与这个切向正交，等价于属于两个法向的线性张成空间。若约束梯度平行或有一个为零，乘子方程的常规必要性就失去保证，应把这些退化点另列候选。

消元时还会人为丢失分支。例如得到$\lambda(x-y)=0$，正确分支是$\lambda=0$或$x=y$；未经讨论就约去$\lambda$，会漏掉原本满足方程的点。例~\ref{ex:c3-e7}中的全局最小值正来自这个零乘子分支。平方、开方及参数化同样要跟踪等号和符号条件。

闭且有界的非空约束集合上，连续目标函数有全局极值，但乘子法只负责正规内点的候选。带非负约束的熵问题还包含各条边与顶点；包围几何体的问题则包含一整段连续不等式，单个接触条件只是必要信息。下面的例题分别保留这些结构，使“求驻点”与“证明全局最优”成为相接的两步。

\begin{example}[三个二元函数的驻点分类]\label{ex:c3-e3}
讨论$f=3x^3-12xy+y^3$、$g=\ln(1+x^2+y^2)$、$h=\ee^{x^2+xy+y^2}$的全部驻点及性质。 见图~\ref{fig:c3-e3-q}。
\end{example}

\begin{figure}[H]\centering
\begin{tikzpicture}[>=Stealth,line cap=round,line join=round]
\begin{scope}[x=0.85000cm,y=0.85000cm]\draw[axis](-.3,0)--(3,0)node[right]{$x$};\draw[axis](0,-.3)--(0,4)node[above]{$y$};\draw[curve,domain=0:2.2]plot(\x,{.75*(\x)^2});\draw[dashed,domain=0:3.4]plot({.25*(\x)^2},\x);\fill(0,0)circle(1.5pt)node[below left]{$O$};\fill(1.922999,2.773445)circle(1.5pt)node[right]{$P$};\node[right]at(3.4,2.5){实线：$f_x=0$};\node[right]at(3.4,1.7){虚线：$f_y=0$};\end{scope}
\end{tikzpicture}
\caption{三次函数的两条驻点条件相交于$O$和$P$，其余两函数另由梯度判断。}\label{fig:c3-e3-q}
\end{figure}
\begin{solution}
驻点计算与全局极值判断应分开。三个函数都有全平面定义域，但第一式是可向无穷下降的三次多项式，后两式则可以由非负二次型直接限制下界。图~\ref{fig:c3-e3-q}先用驻点方程的交点找候选，再按Hessian或全局不等式分类。


第一式的偏导为 $f_x=9x^2-12y$、$f_y=-12x+3y^2$。令两者为零，由 $y=3x^2/4$ 代入 $y^2=4x$，得 $9x^4/16=4x$，即 $x(9x^3-64)=0$。所以驻点为 $O$ 和
\[P=\left(\frac4{3^{2/3}},\frac4{3^{1/3}}\right).\]
Hessian矩阵是$\begin{pmatrix}18x&-12\\-12&6y\end{pmatrix}$。在$O$处行列式$-144<0$，为鞍点；在$P$处$xy=16/3$，行列式$108xy-144=432>0$且$18x>0$，为严格局部极小值，其值为$-64/3$。沿$x\to-\infty,y=0$，$f\to-\infty$，因此$f$在整个平面上没有下界，$P$只给出局部极小值。

第二式的梯度为$2(x,y)/(1+x^2+y^2)$，只有原点驻点。由于$g\ge0$且等号仅在原点成立，原点是唯一全局最小值点。第三式指数始终正，驻点方程为$2x+y=0,x+2y=0$，仍只有原点；对$y$求偏导时，$y^2$贡献$2y$。又
\[x^2+xy+y^2=(x+y/2)^2+3y^2/4\ge0,\]
故$h$的全局最小值为$1$。图~\ref{fig:c3-e3-s}的双曲线型与封闭椭圆型分别对应鞍点和正定二次型。


消元关系 $y=3x^2/4$、$y^2=4x$ 还保证两个坐标均非负，没有遗漏负的实驻点。对非零点，$x^3=64/9,y^3=64/3,xy=16/3$，代入目标得 $64/3-64+64/3=-64/3$。

原点的二次主部为$-12xy$，沿$y=x$为负、沿$y=-x$为正，所以图~\ref{fig:c3-e3-s}中的双曲线型等值线对应鞍点。后两个函数在原点均取得全局下界；对于指数二次型，正交主方向的系数为$3/2$与$1/2$，等值线为封闭椭圆。局部二阶正定与全局代数下界在这两题中互相印证，但第一式只有局部极小。

\begin{figure}[H]\centering
\begin{tikzpicture}[>=Stealth,line cap=round,line join=round]
\draw[axis](-1.5,0)--(1.5,0);\draw[axis](0,-1.4)--(0,1.5);\draw[curve,domain=0.24:1.3,samples=61]plot(\x,{0.3/\x});\draw[curve,domain=-1.3:-0.24,samples=61]plot(\x,{0.3/\x});\draw[curve,domain=0.24:1.3,samples=61]plot(\x,{-0.3/\x});\draw[curve,domain=-1.3:-0.24,samples=61]plot(\x,{-0.3/\x});\node at(0,-1.9){$O$附近：$xy=\text{常数}$};\begin{scope}[shift={(5,0)},rotate=-45]\draw[curve](0,0)ellipse(1.2 and .69282);\draw[curve](0,0)ellipse(.7 and .40415);\fill(0,0)circle(1.5pt);\end{scope}\node at(5,-1.9){$x^2+xy+y^2$的等值线};
\end{tikzpicture}
\caption{左图为鞍点处二次主项的双曲线，右图为正定二次型的椭圆。}\label{fig:c3-e3-s}
\end{figure}
\end{solution}

\Needspace{8\baselineskip}


\Needspace{8\baselineskip}


\Needspace{8\baselineskip}


\Needspace{8\baselineskip}

\subsection{势能的固定变量与允许位移}
对势能$U(\vct r)$，静力平衡要求沿所有允许位移的一次变化为零。自由质点可以沿任意方向移动，故需$\nabla U=0$；受约束质点只沿切空间运动，法向梯度可由约束力平衡。求二阶变化时也应保持同一组条件。第一讲圆锥小振动固定轴向角动量，近圆轨道固定中心角动量，所求的是相应有效势的曲率。

若局部允许坐标为$\vct q$，二阶动能和势能写成
\[T_2=\frac12\dot{\vct q}^{\mathsf T}M\dot{\vct q},\qquad
U_2=\frac12\vct q^{\mathsf T}K\vct q,\]
则线性方程是$M\ddot{\vct q}+K\vct q=0$，频率满足$K\vct u=\omega^2M\vct u$。这里$M$由允许运动的几何决定，$K$由保持守恒量或约束后的势能决定。只有二者使用同一坐标，频率才有确定意义。下面先讨论自由质点的二维势阱，随后讨论直导轨和曲面约束。
\begin{exercise}[各向异性势阱中的平衡]\label{ex:c3-h7}
质点势能$U=\frac12k_xx^2+cxy+\frac12k_yy^2$，$k_x,k_y>0$。求力场与平衡稳定条件，并在$k_x=k_y=k$时求正交主方向。 见图~\ref{fig:c3-h7-q}。
\end{exercise}

\begin{figure}[H]\centering
\begin{tikzpicture}[>=Stealth,line cap=round,line join=round]
\draw[axis](-1.7,0)--(1.9,0)node[right]{$x$};\draw[axis](0,-1.6)--(0,1.9)node[above]{$y$};
\draw[curve,rotate=45](0,0)ellipse(.8 and 1.38564);
\node[right,align=left]at(2.3,.9){$k_x=k_y=k$\\[4pt]$c=k/2$};
\end{tikzpicture}
\caption{固定图示参数后，等势椭圆长轴沿$x=-y$，与较小主刚度$k-c$对应。}\label{fig:c3-h7-q}
\end{figure}
\begin{solution}
势能是二次型，力场为其负梯度。稳定性由两个主刚度共同决定，不能只检查$k_x,k_y$分别为正，因为交叉项可形成斜向下降方向。图~\ref{fig:c3-h7-q}给出一组明确的正定参数，使等势椭圆主轴与坐标轴不同。


$\vct F=-\nabla U=(-k_xx-cy,-cx-k_yy)$。势阱严格稳定要求Hessian正定，即
\[\boxed{k_xk_y-c^2>0}.
\]
若小于零，有一个下降方向，平衡不稳定；等于零时存在零刚度方向，纯二次模型不能提供严格束缚。图~\ref{fig:c3-h7-s}在$k_x=k_y=k$时取$\xi=(x+y)/\sqrt2,\eta=(x-y)/\sqrt2$，得到
\[U=\tfrac12(k+c)\xi^2+\tfrac12(k-c)\eta^2.\]
两个主方向是两条角平分线，稳定条件为$|c|<k$。若质量为$m$，主频为$\sqrt{(k+c)/m}$和$\sqrt{(k-c)/m}$。当$|c|\to k$，一个主频趋零，表示沿相应方向恢复力消失。


对一般 $k_x,k_y$，刚度矩阵为 $K=\begin{pmatrix}k_x&c\\c&k_y\end{pmatrix}$。特征方程 $\det(K-qI)=(k_x-q)(k_y-q)-c^2=0$ 展开为 $q^2-(k_x+k_y)q+k_xk_y-c^2=0$，用二次方程求根得
\[k_\pm=\frac{k_x+k_y\pm\sqrt{(k_x-k_y)^2+4c^2}}2.\]
迹恒正，所以两主刚度同时为正等价于行列式为正。若行列式为零，平衡点构成一条零势能直线；沿该方向的任意小初速度会产生无界匀速漂移，孤立位置在完整位置--速度状态空间中并不稳定。若行列式为负，负刚度方向则出现指数增长。解答图取$k_x=k_y=k,c=k/2$，短轴沿$\xi$、长轴沿$\eta$，与两刚度$3k/2,k/2$一一对应。

一般情形可先配方：
\[U=\frac{k_x}{2}\left(x+\frac c{k_x}y\right)^2
+\frac{k_xk_y-c^2}{2k_x}y^2.\]
由于$k_x>0$，第二项系数决定是否在所有非零方向升高。当行列式为负，选择$x=-cy/k_x$即可得到负势能方向；为零时这整条直线势能相同。

$k_x=k_y=k$时，图~\ref{fig:c3-h7-s}的正交主坐标把动能也保留为$m(\dot\xi^2+\dot\eta^2)/2$，所以两条独立方程分别为$m\ddot\xi+(k+c)\xi=0$与$m\ddot\eta+(k-c)\eta=0$。主频需要同时使用主刚度和同一质量，而不是仅从椭圆半轴读出频率数值。

\begin{figure}[H]\centering
\begin{tikzpicture}[>=Stealth,line cap=round,line join=round]
\draw[axis](-1.5,-1.5)--(1.5,1.5)node[above]{$\xi$};
\draw[axis](-1.5,1.5)--(1.5,-1.5)node[right]{$\eta$};
\draw[curve,rotate=45](0,0)ellipse(.8 and 1.38564);
\node[align=left,right]at(2.3,.7){$\xi$方向：$k+c$\\[5pt]$\eta$方向：$k-c$};
\end{tikzpicture}
\caption{主坐标与题图采用同一组参数$c=k/2$；两方向的半轴比为$1:\sqrt3$，与刚度平方根成反比。}\label{fig:c3-h7-s}
\end{figure}
\end{solution}
\begin{exercise}[导轨能否约束住一个鞍点]\label{ex:c3-v1}

平面内质点的势能为$U=k(x^2+4xy+y^2)/2$，$k>0$，质量$m$。先判断原点对自由质点的稳定性。再把小环限制在固定光滑直导轨$y=\lambda x$上，求原点稳定的$\lambda$范围及小振动频率。导轨为双侧约束。 见图~\ref{fig:c3-v1-q}。
\end{exercise}

\begin{figure}[H]\centering
\begin{tikzpicture}[>=Stealth,line cap=round,line join=round]

\draw[axis](-1.6,0)--(1.8,0)node[right]{$x$};\draw[axis](0,-1.6)--(0,1.8)node[above]{$y$};
\draw[thick](-1.35,-1.35)--(1.35,1.35)node[above right]{$\lambda=1$};
\draw[dashed](-1.35,1.35)--(1.35,-1.35)node[below right]{$\lambda=-1$};
\fill(0,0)circle(2pt);\node at(-.4,.15){$O$};\draw[fill=white](.7,.7)circle(.1);

\end{tikzpicture}
\caption{同一势能在两条允许方向上具有相反的曲率，限制运动方向会改变稳定性。}\label{fig:c3-v1-q}
\end{figure}
\begin{solution}
自由质点的 Hessian 为 $k\begin{pmatrix}1&2\\2&1\end{pmatrix}$，特征方程为 $(k-q)^2-4k^2=0$，解得 $q=3k,-k$，所以原点是鞍点。图~\ref{fig:c3-v1-q}的 $y=x$ 和 $y=-x$ 正是这两个主方向。

有导轨时只能取$(\delta x,\delta y)=(\delta x,\lambda\delta x)$。直接代入动能与势能：
\[T=\frac12m(1+\lambda^2)\dot x^2,\qquad
U=\frac12 k(1+4\lambda+\lambda^2)x^2.\]
沿轨道对 $x$ 求势能导数，$\dd U/\dd x=k(1+4\lambda+\lambda^2)x$；动能对应的惯性系数为 $m(1+\lambda^2)$，所以运动方程为
\[m(1+\lambda^2)\ddot x+k(1+4\lambda+\lambda^2)x=0.\]
原点有稳定小振动当且仅当 $1+4\lambda+\lambda^2=(\lambda+2)^2-3>0$，即
\[\boxed{\lambda<-2-\sqrt3\quad\text{或}\quad\lambda>-2+\sqrt3},\qquad
\boxed{\omega^2=\frac{k(1+4\lambda+\lambda^2)}{m(1+\lambda^2)}}.\]
在两个边界值，轨道上的势能恒为零，只有自由漂移；初位置扰动而速度为零时可以静止，但任意小的非零切向初速度会累积位移，因而完整位置--速度状态没有稳定束缚。其余中间区间曲率为负，存在指数增长。

图~\ref{fig:c3-v1-s}以弧长$s=\sqrt{1+\lambda^2}\,x$作横轴。对$\lambda=1$，有$U=3ks^2/2$、$\omega^2=3k/m$；对$\lambda=-1$，有$U=-ks^2/2$，不稳定。用$s$时惯性系数为$m$，用$x$时则为$m(1+\lambda^2)$，频率保持一致。

第一讲的圆锥与近圆轨道在固定角动量后构造有效势，本题则由导轨限制可行方向。三者均需在允许的坐标中计算惯性系数和势能曲率。自由空间的Hessian描述任意方向的变化，不能直接代替导轨上的二阶势能。
\begin{figure}[H]\centering
\begin{tikzpicture}[>=Stealth,line cap=round,line join=round]

\begin{scope}[x=2cm,y=1.2cm]
\draw[axis](-1.2,0)--(1.4,0)node[right]{$s/L$};\draw[axis](0,-.8)--(0,1.8)node[above]{$U/(kL^2)$};
\draw[curve,domain=-1:1,samples=81]plot(\x,{1.5*(\x)^2});\draw[dashed,domain=-1:1,samples=81]plot(\x,{-.5*(\x)^2});
\node[right,align=left]at(1.65,.6){实线：$\lambda=1$\\[5pt]虚线：$\lambda=-1$};
\end{scope}

\end{tikzpicture}
\caption{沿弧长比较两个受约束势能；$L$为任意选定的长度标度。}\label{fig:c3-v1-s}
\end{figure}
\end{solution}


\subsection{几何约束与全局极值}


\begin{example}[球面与平面的交线]\label{ex:c3-e7}
在$x^2+y^2+z^2=11,x+y+z=3$上求$T=20+2x+2y+z^2$的全局最大最小值及其取得点。 见图~\ref{fig:c3-e7-q}。
\end{example}

\begin{figure}[H]\centering
\begin{tikzpicture}[>=Stealth,line cap=round,line join=round]
\draw[curve](0,0)circle(2);\draw[curve,rotate=-20](0,.7)ellipse(1.75 and .42);\draw[axis](0,-2.3)--(0,2.7)node[above]{$z$};\fill(0,0)circle(1.4pt)node[left]{$O$};\node[right]at(2.4,.9){平面不经过球心};
\end{tikzpicture}
\caption{平面与球面交成闭圆，交圆的中心偏离球心。}\label{fig:c3-e7-q}
\end{figure}
\begin{solution}
两个约束把可行域限制在一条闭圆上。目标只需$x+y$与$z$，可先用平面关系消去$x+y$，再用球面给出$z$的允许范围。图~\ref{fig:c3-e7-q}说明平面不经过球心，因此高度范围应围绕交圆中心计算，而非直接取整个球的最高、最低点。


先用平面约束消去$x+y=3-z$，得
\[T=26-2z+z^2=25+(z-1)^2.\]
再由$(x-y)^2=2(x^2+y^2)-(x+y)^2$，得到
\[(x-y)^2=22-2z^2-(3-z)^2=16-3(z-1)^2\ge0.
\]
因此图~\ref{fig:c3-e7-s}中的完整可行区间为$|z-1|\le4/\sqrt3$。于是
\[\boxed{T_{\min}=25,\quad (x,y,z)=(3,-1,1),(-1,3,1)},\]
\[\boxed{T_{\max}=91/3,\quad
(x,y,z)=(1\mp2/\sqrt3,1\mp2/\sqrt3,1\pm4/\sqrt3)}.
\]
在最小点$z=1$时，$x+y=2,xy=-3$，给出上述两个点；最大值的区间端点满足$x=y$。

若用乘子$\nabla T=\lambda\nabla(x^2+y^2+z^2)+\mu\nabla(x+y+z)$，则$2=2\lambda x+\mu$、$2=2\lambda y+\mu$。两式相减仅能得到$\lambda(x-y)=0$。$x=y$给出最大点，$\lambda=0$则由$2=\mu,2z=\mu$得到$z=1$，正好给出最小点。若直接除以$\lambda$，会把两个最小点全部漏掉。


可行区间也可由交圆的几何得到。平面法向为$(1,1,1)$，球心到平面的垂足为$C=(1,1,1)$，距离为$\sqrt3$；故交圆半径为$\sqrt{11-3}=2\sqrt2$。竖直单位向量在平面内的投影长度为$\sqrt{2/3}$，因此圆上$z$相对圆心高度的最大偏离为$2\sqrt2\sqrt{2/3}=4/\sqrt3$，与代数范围一致。

图~\ref{fig:c3-e7-s}中两个端点均有$T=91/3$，中点$z=1$有$T=25$。同一个内部高度通常对应交圆上的两个点，因此求得极值高度后，还须代回两条约束确定全部$x,y$取得点。

对任意可行高度$z$，由和与差可将其余坐标完整写成
\[x=\frac{3-z\pm\sqrt{16-3(z-1)^2}}2,\qquad
 y=\frac{3-z\mp\sqrt{16-3(z-1)^2}}2.\]
根号非负不仅是必要条件，也保证构造出实数$x,y$满足两约束，所以区间$|z-1|\le4/\sqrt3$确为完整可达范围。

乘子零分支与代数平方根的两支在此对应，消元时保留它们才能列全取得点。

\begin{figure}[H]\centering
\begin{tikzpicture}[>=Stealth,line cap=round,line join=round]
\begin{scope}[x=.95cm,y=.65cm]
\fill[gray!10]({1-4/sqrt(3)},0)rectangle({1+4/sqrt(3)},{19/3});
\draw[axis](-1.65,0)--(3.7,0)node[right]{$z$};\draw[axis](0,0)--(0,7)node[above]{$T-24$};
\draw[curve,domain=-1.309401:3.309401,samples=121]plot(\x,{(\x-1)^2+1});
\draw[guide]({1-4/sqrt(3)},0)--({1-4/sqrt(3)},{19/3});
\draw[guide]({1+4/sqrt(3)},0)--({1+4/sqrt(3)},{19/3});
\fill(1,1)circle(1.5pt)node[above]{最小};
\fill({1-4/sqrt(3)},{19/3})circle(1.5pt);\fill({1+4/sqrt(3)},{19/3})circle(1.5pt);
\node[above]at({1-4/sqrt(3)},{19/3}){最大};\node[above]at({1+4/sqrt(3)},{19/3}){最大};
\node[below]at({1-4/sqrt(3)},0){$1-4/\sqrt3$};
\node[below]at({1+4/sqrt(3)},0){$1+4/\sqrt3$};
\node[below]at(1,0){$1$};\end{scope}
\end{tikzpicture}
\caption{阴影标明完整可行高度范围；只在此闭区间上比较二次函数，端点与内部驻点都须保留。}\label{fig:c3-e7-s}
\end{figure}
\end{solution}

\Needspace{8\baselineskip}
\begin{example}[包围偏心球的最小旋转椭球]\label{ex:c3-e8}
用中心固定在原点、对称轴为$x$轴的旋转椭球$x^2/a^2+(y^2+z^2)/b^2\le1$包住球$x^2+y^2+z^2\le2x$。求最小体积及$a,b$。长度采用同一单位。 见图~\ref{fig:c3-e8-q}。
\end{example}

\begin{figure}[H]\centering
\begin{tikzpicture}[>=Stealth,line cap=round,line join=round]
\draw[axis](-2.6,0)--(3,0)node[right]{$x$};\draw[axis](0,-1.6)--(0,1.8)node[above]{$y$};\draw[curve](0,0)ellipse(2.3094 and 1.1547);\draw[dashed](1,0)circle(1);\fill(1,0)circle(1.5pt)node[below]{$1$};
\end{tikzpicture}
\caption{取子午面$z=0$，上半平面中$y=\varrho$；绕$x$轴旋转恢复三维模型。}\label{fig:c3-e8-q}
\end{figure}
\begin{solution}
要包住的是整个偏心球，而不是只经过某一个球面点。图~\ref{fig:c3-e8-q}的旋转对称性允许把三维问题降到子午截面的包含关系：对每个$x$，球截面允许的最大横向半径都不能超出椭球截面。这样先得到一整段区间不等式，再优化半轴。


记$A=a^2,B=b^2$，球在截面$x\in[0,2]$处的最大径向距离满足$\varrho^2=2x-x^2$。包围条件等价于
\[g(x)=x^2/A+(2x-x^2)/B\le1\quad(0\le x\le2).
\]
必须控制整个区间。在 $x=1$ 处，$g(1)=1/A+1/B\le1$ 要求 $B>1$；端点 $x=2$ 给出 $A\ge4$。若 $B\ge2$，于是 $\sqrt A B\ge4$，不会优于下文的候选解。

现在取 $1<B<2$。由 $A\ge4>2B$，有 $A>2B$，所以 $g$ 的二阶导数 $g''=2(1/A-1/B)<0$。一阶导数为 $g'=2/B+2x(1/A-1/B)$，令其为零，顶点为
\[x_0=\frac A{A-B}\in(1,2),\qquad
 g(x_0)=\frac A{B(A-B)}.
\]
顶点值不超过 $1$，即 $A\le B(A-B)$，移项为 $(B-1)A\ge B^2$，因此包围条件为 $A\ge B^2/(B-1)$。固定 $B$ 时，体积随 $A$ 增加，最优必取等号，
\[\frac{3V}{4\pi}=\sqrt A B=\frac{B^2}{\sqrt{B-1}}.
\]
记 $W(B)=B^2/\sqrt{B-1}$，则 $\ln W=2\ln B-\tfrac12\ln(B-1)$。求导并通分，
\[
\frac{W'}W=\frac2B-\frac1{2(B-1)}=\frac{3B-4}{2B(B-1)}.
\]
分母为正，唯一零点为 $B=4/3$，导数先负后正。此处 $A=B^2/(B-1)=16/3$，且 $W=16/(3\sqrt3)<4$。结合 $B\ge2$ 的比较，得到全局解
\[\boxed{a=4/\sqrt3,\quad b=2/\sqrt3,\quad
V_{\min}=\frac{64\pi}{9\sqrt3}}.\]
图~\ref{fig:c3-e8-s}中的接触圆位于$x_0=4/3$、$\varrho_0=2\sqrt2/3$。在该处球与椭球的法向平行；把最优$A,B$代回，还能得到
\[1-g(x)=\frac9{16}(x-4/3)^2\ge0,
\]
这给出了整个球确实被包住的直接证明。若原球半径为$R$、中心为$(R,0,0)$，各半轴乘$R$，体积乘$R^3$。


图~\ref{fig:c3-e8-s}中的接触点对应控制最大值的顶点；非负平方式确认其余所有球面截面也被包含。局部相切条件须与这条全区间包含证明一起使用。

\begin{figure}[H]\centering
\begin{tikzpicture}[>=Stealth,line cap=round,line join=round]
\draw[axis](-2.6,0)--(3,0)node[right]{$x$};\draw[axis](0,-1.5)--(0,1.7)node[above]{$y$};\draw[curve](0,0)ellipse(2.309401 and 1.154701);\draw[dashed](1,0)circle(1);\fill(1.333333,.942809)circle(1.7pt)node[above right]{$P$};\draw[guide](1.333333,0)--(1.333333,.942809);\node[below]at(1.333333,0){$4/3$};\draw[axis](1.333333,.942809)--(1.6,1.697);
\end{tikzpicture}
\caption{接触位置由整段包围条件确定，法向平行仅是必要的局部条件。}\label{fig:c3-e8-s}
\end{figure}
\end{solution}

\Needspace{8\baselineskip}
\begin{example}[斜抛物体的最大视角]\label{ex:c3-e9}
物体从原点以速率$v_0$、仰角$45^\circ$抛出，观察者在$(-L,0)$，$L>0$。忽略空气阻力，求物体在飞行期间相对观察者的最大仰角及对应时刻。分别用消元和乘子法说明。 见图~\ref{fig:c3-e9-q}。
\end{example}

\begin{figure}[H]\centering
\begin{tikzpicture}[>=Stealth,line cap=round,line join=round]
\draw[axis](-1.7,0)--(4,0)node[right]{$x$};\draw[axis](0,-.2)--(0,2)node[above]{$y$};\draw[curve,domain=0:3]plot(\x,{\x-(\x)^2/3});\fill(-1,0)circle(1.5pt)node[below]{$A(-L,0)$};\draw[dashed](-1,0)--(1.15,.709);\fill(1.15,.709)circle(1.5pt)node[above]{$P$};\draw[axis](0,0)--(.6,.6)node[above]{$v_0$};
\end{tikzpicture}
\caption{观察者的视线连接固定点$A$与轨迹上的动点$P$。}\label{fig:c3-e9-q}
\end{figure}
\begin{solution}
观察者的仰角由视线斜率$y/(x+L)$决定，物体的最高点不一定给出最大仰角。图~\ref{fig:c3-e9-q}中观察者位于发射点左侧，水平距离随飞行增加，所以高度增加与距离增加对视角的作用需要同时计算。


令$k=g/v_0^2$。轨迹为$x=v_0t/\sqrt2,y=x-kx^2$，飞行区间$0\le x\le1/k$。因为$x+L>0$，视角$\alpha\in[0,\pi/2)$，最大化$\alpha$等价于最大化
\[u(x)=\tan\alpha=\frac{x-kx^2}{x+L}.\]
用商法则求导，
\[
u'=\frac{(1-2kx)(x+L)-(x-kx^2)}{(x+L)^2}
=\frac{L-kx^2-2kLx}{(x+L)^2}.
\]
记分子为 $N(x)$，则 $N'=-2k(x+L)<0$，而 $N(0)=L>0$、$N(1/(2k))=-1/(4k)<0$。因此唯一驻点位于最高点之前，并使导数由正变负。解 $kx^2+2kLx-L=0$，两个根是 $-L\pm\sqrt{L^2+L/k}$，只有加号根在飞行区间内，得到
\[\boxed{x_* =\sqrt{L^2+L/k}-L,\quad t_* =\frac{\sqrt2}{v_0}x_*}.\]
由驻点关系$y_*=(x_*+L)(1-2kx_*)$，
\[\boxed{\alpha_{\max}=\arctan(1-2kx_*)}.
\]
根满足$0<x_*<1/(2k)$，故最大视角在最高点之前取得。

若直接令$f(x,y)=y/(x+L)$、约束$G=y-x+kx^2=0$，乘子方程为
\[-y/(x+L)^2=\lambda(-1+2kx),\qquad 1/(x+L)=\lambda.
\]
消去乘子又得到$y=(x+L)(1-2kx)$。图~\ref{fig:c3-e9-s}中，视线斜率$y/(x+L)$等于轨迹切线斜率$1-2kx$，所以最大视角的视线与轨迹相切。飞行两端的仰角为零，内部驻点为正，故它给出全飞行段的最大值。$L\to\infty$时最优点趋近最高点，$L\to0^+$时最优点趋近发射点，最大仰角趋于$45^\circ$；在$L=0$的退化问题中该值是$t\to0^+$的上确界。



\begin{figure}[H]\centering
\begin{tikzpicture}[>=Stealth,line cap=round,line join=round]
\draw[axis](-1.3,0)--(3.5,0)node[right]{$x$};\draw[axis](0,-.2)--(0,1.8)node[above]{$y$};\draw[curve,domain=0:3]plot(\x,{\x-(\x)^2/3});\draw[thick](-1,0)--(2.8,1.266667);\fill(1,.666667)circle(1.5pt)node[above left]{$P_*$};\draw[guide](1,0)--(1,.666667);\node[below]at(1,0){$x_*$};\node at(2.2,1.65){$L=1,\ k=1/3$};
\end{tikzpicture}
\caption{数值示意中切点$x_*=1$，视线斜率与轨迹斜率均为$1/3$。}\label{fig:c3-e9-s}
\end{figure}
\end{solution}


\begin{exercise}[水平恒力作用下的抛物线导轨]\label{ex:c3-v2}

质量$m$的小环套在竖直平面内的固定光滑导轨$y=x^2/a$上，$a>0$，$y$竖直向上。除重力外，还受水平向右恒力$F\ge0$。求平衡点及其附近的小振动频率，说明直接用$U_{xx}/m$是否正确。 见图~\ref{fig:c3-v2-q}。
\end{exercise}

\begin{figure}[H]\centering
\begin{tikzpicture}[>=Stealth,line cap=round,line join=round]

\begin{scope}[x=2cm,y=1.8cm]
\draw[axis](-1.1,0)--(1.5,0)node[right]{$x/a$};\draw[axis](0,-.1)--(0,1.5)node[above]{$y/a$};
\draw[curve,domain=-1:1.15,samples=81]plot(\x,{(\x)^2});\fill(.5,.25)circle(2pt)node[above left]{$P$};
\draw[axis](.5,.25)--(1.05,.25)node[right]{$F$};\draw[axis](.5,.25)--(.5,-.35)node[below]{$mg$};
\end{scope}

\end{tikzpicture}
\caption{导轨曲线按题设精确绘制；图取$F=mg$时的平衡点$P$。}\label{fig:c3-v2-q}
\end{figure}
\begin{solution}
图~\ref{fig:c3-v2-q}中重力和水平恒力均为保守力，平面势能为$U=mgy-Fx$。沿允许路径$y=x^2/a$代入，得到
\[U_c(x)=\frac{mg}{a}x^2-Fx
=\frac{mg}{a}\left(x-\frac{Fa}{2mg}\right)^2-\frac{F^2a}{4mg}.\]
故唯一平衡点为
\[\boxed{x_0=\frac{Fa}{2mg},\qquad y_0=\frac{F^2a}{4m^2g^2}}.\]
它是沿导轨的全局势能最小点。平面中的$U_{xx}=0$，不能据此判断没有回复力，因为可行的二阶位移还会改变$y$。

由导轨斜率$\dot y=(2x/a)\dot x$，动能是
\[T=\frac12m\left(1+\frac{4x^2}{a^2}\right)\dot x^2.
\]
记 $M(x)=m(1+4x^2/a^2)$。由 $\mathcal L=M(x)\dot x^2/2-U_c(x)$，乘积求导为
\[
\frac{\dd}{\dd t}\frac{\partial\mathcal L}{\partial\dot x}
=\frac{\dd}{\dd t}(M\dot x)=M\ddot x+M'\dot x^2,\qquad
\frac{\partial\mathcal L}{\partial x}=\frac12M'\dot x^2-U_c'.
\]
所以完整方程为 $M\ddot x+\tfrac12M'\dot x^2+U_c'=0$，其中 $M'=8mx/a^2$、$U_c'=2mgx/a-F$。令 $\xi=x-x_0$，$\dot x^2$ 项为二阶小量，$M$ 的变化与 $\ddot\xi$ 的乘积也为二阶；线性项只取 $M(x_0)$，而 $U_c''=2mg/a$，故
\[m\left[1+\left(\frac F{mg}\right)^2\right]\ddot\xi
+\frac{2mg}{a}\xi=0,\qquad
\boxed{\omega^2=\frac{2g/a}{1+(F/mg)^2}}.\]
即便使用沿约束代入后的$U_c''$，若只除以$m$仍会漏掉动能的几何因子。

图~\ref{fig:c3-v2-s}显示平衡点切向增量满足$\delta y=(2x_0/a)\delta x$。弧长微扰为$\delta s=\sqrt{1+(F/mg)^2}\,\delta x$，将二次势能也改写成$\delta s$，其系数恰好是$m\omega^2$。水平力增大把平衡点移向更陡的轨段，同一水平位移对应更长的实际路程，因而频率减小。$F=0$时恢复$\omega=\sqrt{2g/a}$，与导轨底部曲率半径$a/2$的结果一致。
\begin{figure}[H]\centering
\begin{tikzpicture}[>=Stealth,line cap=round,line join=round]

\begin{scope}[x=3cm,y=3cm]
\draw[axis](-.15,0)--(1.25,0)node[right]{$x/a$};\draw[axis](0,-.1)--(0,1.15)node[above]{$y/a$};
\draw[curve,domain=0:1.05,samples=81]plot(\x,{(\x)^2});
\draw[dashed](.2,-.05)--(1, .75);\fill(.5,.25)circle(1.5pt)node[above left]{$P$};
\draw[guide](.5,.25)--(.85,.25)--(.85,.6);
\draw[axis](.5,.25)--(.85,.6)node[midway,above left]{$\delta s$};
\node[below]at(.675,.25){$\delta x$};\node[right]at(.85,.425){$\delta y$};
\end{scope}

\end{tikzpicture}
\caption{取$F=mg$，在真实导轨上用切线构造一阶位移三角形；虚线为$P$点切线。}\label{fig:c3-v2-s}
\end{figure}
\end{solution}


\subsection{数学应用：竞争、凹性与熵}
以下模型讨论单边偏离、严格凹性和带非负边界的约束。驻点和乘子方程给出候选，极值是否全局成立还取决于可行集合与边界。
\begin{example}[竞争、合作与单边偏离]\label{ex:c3-e4}
两人的非负决策量为$x,y$，收益为$I_1=(4-x-y)x,I_2=(4-x-y)y$。求相互最优反应的交点、双方取相同决策时的合作最优值，以及对方坚持合作时单方最优偏离。若偏离以概率$0<\beta\le1$被发现并扣除$\Delta I$，求阻止偏离的最低扣除量。再考虑$n$家企业，价格为$6-\sum_jx_j$，单位成本为$2$，相应收益为$I_i=(4-\sum_jx_j)x_i$。求其相互最优反应的交点。 见图~\ref{fig:c3-e4-q}。
\end{example}

\begin{figure}[H]\centering
\begin{tikzpicture}[>=Stealth,line cap=round,line join=round]
\draw[axis](0,0)--(4.5,0)node[right]{$x$};\draw[axis](0,0)--(0,4.5)node[above]{$y$};\draw[curve](0,4)--(2,0)node[below]{$2$};\draw[dashed](0,2)--(4,0);\fill(1.3333,1.3333)circle(1.5pt)node[above right]{$N$};
\end{tikzpicture}
\caption{两条最优反应直线相交，但合作点不在交点。}\label{fig:c3-e4-q}
\end{figure}
\begin{solution}
竞争、合作与单边偏离对应三个不同的最优化条件。固定对手决策求自己的最优，是最优反应；双方共同选择以提高收益，是合作；合作点能否抵抗单方改变，还要再解一个固定对手的最优化。图~\ref{fig:c3-e4-q}中的两条反应线只决定第一种问题。


固定 $y$，$I_1=(4-y)x-x^2$ 的导数为 $4-y-2x$、二阶导数为 $-2$，所以最优反应为 $x=\max\{0,(4-y)/2\}$。另一人同理。在内点联立 $2x+y=4,x+2y=4$，得到
\[x=y=4/3,\qquad I_1=I_2=16/9.\]
非负边界不会形成额外的相互最优点：若 $x=0$，对方最优反应为 $y=2$，此时第一人又会选择 $x=1$，所以边界不是均衡。若双方约定 $x=y=q$，每人收益为 $4q-2q^2$，求导得 $4-4q=0$，在 $q=1$ 达到 $2$。对方坚持 $y=1$ 时，第一人收益为 $3x-x^2$，求导得 $3-2x=0$，最优量为 $x=3/2$，收益 $9/4$，比合作多 $1/4$。

图~\ref{fig:c3-e4-s}把合作点$C=(1,1)$与偏离点$D=(3/2,1)$分开。发现与未发现两种结果下，偏离的期望收益是
\[\beta(9/4-\Delta I)+(1-\beta)9/4=9/4-\beta\Delta I.\]
要使之不超过合作收益$2$，必须且只需
\[\boxed{\Delta I\ge\frac1{4\beta}}.\]
这里采用的是题设的一次性扣除机制；若另有后续收益变化，应另建模型，不能把两种惩罚机制混加。

$n$人情形，对$x_i$求偏导得$4-\sum_jx_j-x_i=0$。内点对称解满足$(n+1)x_i=4$，因此
\[\boxed{x_i=\frac4{n+1},\qquad I_i=\frac{16}{(n+1)^2}}.\]


合作时总收益为$(4-x-y)(x+y)$，只依赖总量$s=x+y$，在$s=2$达到$4$；题目进一步要求相同决策，才选出$x=y=1$并各得$2$。

$n$人推广中，若有$k$家正产量企业，正产量者的内点条件使其产量都为$4/(k+1)$；若$k<n$，零产量者面对总量$4k/(k+1)<4$会有正边际收益，故也要进入。因此$k=n$，对称内点解不仅存在，而且覆盖全部非负均衡。

\begin{figure}[H]\centering
\begin{tikzpicture}[>=Stealth,line cap=round,line join=round]
\draw[axis](0,0)--(3.2,0)node[right]{$x$};\draw[axis](0,0)--(0,2.6)node[above]{$I_1$};\draw[curve,domain=0:3]plot(\x,{(3-\x)*\x});\fill(1,2)circle(1.5pt)node[above left]{$C$};\fill(1.5,2.25)circle(1.5pt)node[above right]{$D$};\draw[guide](1,0)--(1,2);\draw[guide](1.5,0)--(1.5,2.25);\node at(2.4,1.8){$y=1$};
\end{tikzpicture}
\caption{合作收益为$2$，在对方不变时偏离可增加$1/4$。}\label{fig:c3-e4-s}
\end{figure}
\end{solution}

\begin{example}[凹二次函数及线性约束]\label{ex:c3-e5}
求$P=-5x^2-8y^2-2xy+42x+102y$在整个平面及约束$3x+5y=100$上的最大值。 见图~\ref{fig:c3-e5-q}。
\end{example}

\begin{figure}[H]\centering
\begin{tikzpicture}[>=Stealth,line cap=round,line join=round]
\draw[axis](-2.9,0)--(3.5,0)node[right]{$u/5$};\draw[axis](0,-2.5)--(0,3)node[above]{$v/5$};\draw[curve]plot coordinates{(-2.6035561,0.7882934) (-2.5780993,0.8650432) (-2.5486170,0.9404423) (-2.5151554,1.0143730) (-2.4777666,1.0867199) (-2.4365091,1.1573701) (-2.3914473,1.2262131) (-2.3426516,1.2931415) (-2.2901980,1.3580509) (-2.2341686,1.4208398) (-2.1746508,1.4814103) (-2.1117376,1.5396677) (-2.0455271,1.5955211) (-1.9761227,1.6488832) (-1.9036329,1.6996709) (-1.8281708,1.7478047) (-1.7498542,1.7932095) (-1.6688055,1.8358145) (-1.5851511,1.8755530) (-1.4990216,1.9123631) (-1.4105516,1.9461872) (-1.3198792,1.9769726) (-1.2271460,2.0046712) (-1.1324967,2.0292398) (-1.0360792,2.0506399) (-0.9380440,2.0688382) (-0.8385441,2.0838062) (-0.7377349,2.0955207) (-0.6357739,2.1039632) (-0.5328201,2.1091207) (-0.4290344,2.1109849) (-0.3245789,2.1095532) (-0.2196165,2.1048276) (-0.1143112,2.0968156) (-0.0088275,2.0855297) (0.0966700,2.0709874) (0.2020166,2.0532116) (0.3070478,2.0322299) (0.4115996,2.0080751) (0.5155086,1.9807850) (0.6186128,1.9504021) (0.7207511,1.9169738) (0.8217640,1.8805525) (0.9214939,1.8411949) (1.0197849,1.7989624) (1.1164836,1.7539211) (1.2114391,1.7061413) (1.3045031,1.6556975) (1.3955303,1.6026686) (1.4843785,1.5471373) (1.5709090,1.4891903) (1.6549867,1.4289181) (1.7364804,1.3664148) (1.8152628,1.3017781) (1.8912108,1.2351087) (1.9642059,1.1665109) (2.0341342,1.0960917) (2.1008864,1.0239611) (2.1643583,0.9502317) (2.2244509,0.8750187) (2.2810702,0.7984394) (2.3341279,0.7206134) (2.3835411,0.6416623) (2.4292327,0.5617093) (2.4711314,0.4808792) (2.5091717,0.3992983) (2.5432942,0.3170940) (2.5734456,0.2343946) (2.5995790,0.1513291) (2.6216533,0.0680274) (2.6396343,-0.0153805) (2.6534939,-0.0987644) (2.6632103,-0.1819941) (2.6687684,-0.2649396) (2.6701595,-0.3474715) (2.6673815,-0.4294608) (2.6604387,-0.5107796) (2.6493420,-0.5913009) (2.6341086,-0.6708989) (2.6147624,-0.7494493) (2.5913335,-0.8268297) (2.5638586,-0.9029190) (2.5323805,-0.9775985) (2.4969484,-1.0507516) (2.4576176,-1.1222640) (2.4144495,-1.1920242) (2.3675116,-1.2599232) (2.3168770,-1.3258550) (2.2626250,-1.3897166) (2.2048401,-1.4514083) (2.1436126,-1.5108339) (2.0790381,-1.5679004) (2.0112175,-1.6225189) (1.9402566,-1.6746039) (1.8662662,-1.7240743) (1.7893618,-1.7708528) (1.7096636,-1.8148662) (1.6272959,-1.8560460) (1.5423874,-1.8943278) (1.4550706,-1.9296518) (1.3654820,-1.9619629) (1.2737613,-1.9912106) (1.1800517,-2.0173493) (1.0844997,-2.0403381) (0.9872543,-2.0601412) (0.8884675,-2.0767276) (0.7882934,-2.0900714) (0.6868885,-2.1001519) (0.5844110,-2.1069532) (0.4810212,-2.1104647) (0.3768802,-2.1106811) (0.2721508,-2.1076018) (0.1669965,-2.1012318) (0.0615814,-2.0915809) (-0.0439299,-2.0786643) (-0.1493725,-2.0625021) (-0.2545819,-2.0431195) (-0.3593938,-2.0205469) (-0.4636446,-1.9948194) (-0.5671714,-1.9659773) (-0.6698127,-1.9340655) (-0.7714081,-1.8991338) (-0.8717991,-1.8612370) (-0.9708289,-1.8204340) (-1.0683428,-1.7767886) (-1.1641886,-1.7303690) (-1.2582167,-1.6812476) (-1.3502803,-1.6295011) (-1.4402355,-1.5752104) (-1.5279420,-1.5184602) (-1.6132628,-1.4593390) (-1.6960646,-1.3979393) (-1.7762183,-1.3343569) (-1.8535986,-1.2686910) (-1.9280847,-1.2010443) (-1.9995604,-1.1315222) (-2.0679139,-1.0602334) (-2.1330387,-0.9872891) (-2.1948330,-0.9128034) (-2.2532003,-0.8368924) (-2.3080495,-0.7596747) (-2.3592950,-0.6812708) (-2.4068567,-0.6018032) (-2.4506604,-0.5213960) (-2.4906377,-0.4401747) (-2.5267261,-0.3582661) (-2.5588694,-0.2757981) (-2.5870173,-0.1928995) (-2.6111258,-0.1096997) (-2.6311574,-0.0263286) (-2.6470808,0.0570836) (-2.6588710,0.1404067) (-2.6665098,0.2235105) (-2.6699851,0.3062653) (-2.6692915,0.3885420) (-2.6644302,0.4702120) (-2.6554087,0.5511478) (-2.6422410,0.6312231) (-2.6249479,0.7103128) (-2.6035561,0.7882934)};\draw[curve]plot coordinates{(-1.3017781,0.3941467) (-1.2890496,0.4325216) (-1.2743085,0.4702211) (-1.2575777,0.5071865) (-1.2388833,0.5433600) (-1.2182546,0.5786850) (-1.1957237,0.6131065) (-1.1713258,0.6465708) (-1.1450990,0.6790254) (-1.1170843,0.7104199) (-1.0873254,0.7407051) (-1.0558688,0.7698338) (-1.0227635,0.7977605) (-0.9880614,0.8244416) (-0.9518165,0.8498354) (-0.9140854,0.8739024) (-0.8749271,0.8966048) (-0.8344027,0.9179072) (-0.7925755,0.9377765) (-0.7495108,0.9561815) (-0.7052758,0.9730936) (-0.6599396,0.9884863) (-0.6135730,1.0023356) (-0.5662484,1.0146199) (-0.5180396,1.0253199) (-0.4690220,1.0344191) (-0.4192720,1.0419031) (-0.3688675,1.0477603) (-0.3178869,1.0519816) (-0.2664101,1.0545603) (-0.2145172,1.0554925) (-0.1622894,1.0547766) (-0.1098082,1.0524138) (-0.0571556,1.0484078) (-0.0044137,1.0427648) (0.0483350,1.0354937) (0.1010083,1.0266058) (0.1535239,1.0161149) (0.2057998,1.0040376) (0.2577543,0.9903925) (0.3093064,0.9752010) (0.3603756,0.9584869) (0.4108820,0.9402762) (0.4607469,0.9205974) (0.5098924,0.8994812) (0.5582418,0.8769606) (0.6057196,0.8530707) (0.6522516,0.8278488) (0.6977651,0.8013343) (0.7421892,0.7735686) (0.7854545,0.7445951) (0.8274934,0.7144590) (0.8682402,0.6832074) (0.9076314,0.6508890) (0.9456054,0.6175544) (0.9821030,0.5832555) (1.0170671,0.5480459) (1.0504432,0.5119806) (1.0821792,0.4751159) (1.1122254,0.4375093) (1.1405351,0.3992197) (1.1670639,0.3603067) (1.1917706,0.3208311) (1.2146164,0.2808546) (1.2355657,0.2404396) (1.2545858,0.1996492) (1.2716471,0.1585470) (1.2867228,0.1171973) (1.2997895,0.0756646) (1.3108267,0.0340137) (1.3198172,-0.0076902) (1.3267469,-0.0493822) (1.3316051,-0.0909970) (1.3343842,-0.1324698) (1.3350798,-0.1737357) (1.3336908,-0.2147304) (1.3302194,-0.2553898) (1.3246710,-0.2956504) (1.3170543,-0.3354494) (1.3073812,-0.3747247) (1.2956667,-0.4134148) (1.2819293,-0.4514595) (1.2661902,-0.4887992) (1.2484742,-0.5253758) (1.2288088,-0.5611320) (1.2072248,-0.5960121) (1.1837558,-0.6299616) (1.1584385,-0.6629275) (1.1313125,-0.6948583) (1.1024200,-0.7257042) (1.0718063,-0.7554169) (1.0395191,-0.7839502) (1.0056087,-0.8112594) (0.9701283,-0.8373020) (0.9331331,-0.8620372) (0.8946809,-0.8854264) (0.8548318,-0.9074331) (0.8136479,-0.9280230) (0.7711937,-0.9471639) (0.7275353,-0.9648259) (0.6827410,-0.9809814) (0.6368806,-0.9956053) (0.5900259,-1.0086746) (0.5422498,-1.0201690) (0.4936271,-1.0300706) (0.4442337,-1.0383638) (0.3941467,-1.0450357) (0.3434442,-1.0500759) (0.2922055,-1.0534766) (0.2405106,-1.0552324) (0.1884401,-1.0553405) (0.1360754,-1.0538009) (0.0834982,-1.0506159) (0.0307907,-1.0457905) (-0.0219649,-1.0393321) (-0.0746862,-1.0312510) (-0.1272910,-1.0215598) (-0.1796969,-1.0102734) (-0.2318223,-0.9974097) (-0.2835857,-0.9829886) (-0.3349063,-0.9670327) (-0.3857041,-0.9495669) (-0.4358996,-0.9306185) (-0.4854144,-0.9102170) (-0.5341714,-0.8883943) (-0.5820943,-0.8651845) (-0.6291084,-0.8406238) (-0.6751401,-0.8147506) (-0.7201178,-0.7876052) (-0.7639710,-0.7592301) (-0.8066314,-0.7296695) (-0.8480323,-0.6989697) (-0.8881091,-0.6671785) (-0.9267993,-0.6343455) (-0.9640424,-0.6005221) (-0.9997802,-0.5657611) (-1.0339570,-0.5301167) (-1.0665194,-0.4936446) (-1.0974165,-0.4564017) (-1.1266002,-0.4184462) (-1.1540248,-0.3798373) (-1.1796475,-0.3406354) (-1.2034284,-0.3009016) (-1.2253302,-0.2606980) (-1.2453188,-0.2200874) (-1.2633631,-0.1791331) (-1.2794347,-0.1378991) (-1.2935086,-0.0964498) (-1.3055629,-0.0548498) (-1.3155787,-0.0131643) (-1.3235404,0.0285418) (-1.3294355,0.0702033) (-1.3332549,0.1117552) (-1.3349925,0.1531327) (-1.3346458,0.1942710) (-1.3322151,0.2351060) (-1.3277043,0.2755739) (-1.3211205,0.3156115) (-1.3124739,0.3551564) (-1.3017781,0.3941467)};\draw[thick](-.5,2.74)--(3,.64);\fill(0,0)circle(1.5pt)node[below left]{$(3,6)$};\node at(3.7,2.2){$3u+5v=61$};
\end{tikzpicture}
\caption{以无约束最优点为原点，等值线与实际约束直线相切；$u=x-3,v=y-6$。}\label{fig:c3-e5-q}
\end{figure}
\begin{solution}
二次函数的Hessian恒定，先找到无约束驻点，再平移到该点，可把约束最优写成“在直线上找最小正定二次型”。图~\ref{fig:c3-e5-q}中的等值椭圆以$(3,6)$为中心，约束线并不经过中心，因此最大可达等值线在某一点与直线相切。


对原函数求偏导，$P_x=-10x-2y+42$、$P_y=-2x-16y+102$。令两者为零，第一式给 $y=21-5x$，代入第二式得 $2x+16(21-5x)=102$，即 $78x=234$，解得 $(x,y)=(3,6)$、$P=369$。对增量 $u=x-3,v=y-6$，配方为
\[P=369-(5u^2+2uv+8v^2).
\]
括号内正定，故这是唯一全局最大值。

有约束时 $3u+5v=61$。设 $Q=\begin{pmatrix}5&1\\1&8\end{pmatrix}$、$c=(3,5)^{\mathsf T}$，最小化 $\vct u^{\mathsf T}Q\vct u$ 给出 $2Q\vct u=\lambda c$。计算 $\det Q=5\cdot8-1=39$，所以 $Q^{-1}=\tfrac1{39}\begin{pmatrix}8&-1\\-1&5\end{pmatrix}$，$Q^{-1}c=(19,22)^{\mathsf T}/39$、$c^{\mathsf T}Q^{-1}c=(3\cdot19+5\cdot22)/39=167/39$。
由 $\vct u=(\lambda/2)Q^{-1}c$ 再代入约束，$61=(\lambda/2)c^{\mathsf T}Q^{-1}c$，得到 $\lambda/2=61\cdot39/167$，于是
\[\vct u=\frac{61}{167}(19,22)^{\mathsf T},\qquad
\boxed{x=1660/167,\quad y=2344/167}.
\]
在最优点 $Q\vct u=(\lambda/2)c$，故损失量为 $\vct u^{\mathsf T}Q\vct u=(\lambda/2)c^{\mathsf T}\vct u=(61\cdot39/167)\cdot61=145119/167$，因此约束最大值为
\[\boxed{P_{\max}=-83496/167}.\]
图~\ref{fig:c3-e5-s}中的切点具有梯度法向一致的性质。约束把决策点推离原最大点很远，函数值可以变为负；不能因为题目叫“最大值”便预设它为正。二次型严格正定，保证此切点是整条直线上的唯一最优点。


上述最小损失也可从矩阵形式直接写出：
\[\vct u^{\mathsf T}Q\vct u=\frac{61^2}{c^{\mathsf T}Q^{-1}c}.\]
图中的 $P_*$ 是约束直线与最高可达等值线的切点。若用一元消元法，令 $y=(100-3x)/5$ 代入原函数，整理为
\[
P_c(x)=-\frac{167}{25}x^2+\frac{664}{5}x-1160,
\qquad P_c'=-\frac{334}{25}x+\frac{664}{5}.
\]
令 $P_c'=0$，得 $x=1660/167$，再代回约束得 $y=2344/167$；$P_c''=-334/25<0$，与矩阵法一致。

直接的乘子方程为
\[42-10x-2y=3\mu,\qquad102-2x-16y=5\mu,\qquad3x+5y=100.\]
这是一组线性方程，可直接消元得到相同坐标。平移后的矩阵法把计算中的几何意义保留下来：$Q^{-1}c$给出由约束法向经二次型度量修正后的位移方向，通常不与欧氏法向$c$平行。

图~\ref{fig:c3-e5-s}中目标梯度朝向椭圆中心，约束法向与它共线。对沿约束直线的任意非零位移 $\vct h$，有 $c^{\mathsf T}\vct h=0$，交叉项因乘子关系消失，故 $P(P_*+\vct h)=P(P_*)-\vct h^{\mathsf T}Q\vct h<P(P_*)$。这直接证明了沿整条直线的唯一全局最大性。

\begin{figure}[H]\centering
\begin{tikzpicture}[>=Stealth,line cap=round,line join=round]
\draw[axis](-2.9,0)--(3.5,0)node[right]{$u/5$};\draw[axis](0,-2.5)--(0,3)node[above]{$v/5$};\draw[curve]plot coordinates{(-2.6035561,0.7882934) (-2.5780993,0.8650432) (-2.5486170,0.9404423) (-2.5151554,1.0143730) (-2.4777666,1.0867199) (-2.4365091,1.1573701) (-2.3914473,1.2262131) (-2.3426516,1.2931415) (-2.2901980,1.3580509) (-2.2341686,1.4208398) (-2.1746508,1.4814103) (-2.1117376,1.5396677) (-2.0455271,1.5955211) (-1.9761227,1.6488832) (-1.9036329,1.6996709) (-1.8281708,1.7478047) (-1.7498542,1.7932095) (-1.6688055,1.8358145) (-1.5851511,1.8755530) (-1.4990216,1.9123631) (-1.4105516,1.9461872) (-1.3198792,1.9769726) (-1.2271460,2.0046712) (-1.1324967,2.0292398) (-1.0360792,2.0506399) (-0.9380440,2.0688382) (-0.8385441,2.0838062) (-0.7377349,2.0955207) (-0.6357739,2.1039632) (-0.5328201,2.1091207) (-0.4290344,2.1109849) (-0.3245789,2.1095532) (-0.2196165,2.1048276) (-0.1143112,2.0968156) (-0.0088275,2.0855297) (0.0966700,2.0709874) (0.2020166,2.0532116) (0.3070478,2.0322299) (0.4115996,2.0080751) (0.5155086,1.9807850) (0.6186128,1.9504021) (0.7207511,1.9169738) (0.8217640,1.8805525) (0.9214939,1.8411949) (1.0197849,1.7989624) (1.1164836,1.7539211) (1.2114391,1.7061413) (1.3045031,1.6556975) (1.3955303,1.6026686) (1.4843785,1.5471373) (1.5709090,1.4891903) (1.6549867,1.4289181) (1.7364804,1.3664148) (1.8152628,1.3017781) (1.8912108,1.2351087) (1.9642059,1.1665109) (2.0341342,1.0960917) (2.1008864,1.0239611) (2.1643583,0.9502317) (2.2244509,0.8750187) (2.2810702,0.7984394) (2.3341279,0.7206134) (2.3835411,0.6416623) (2.4292327,0.5617093) (2.4711314,0.4808792) (2.5091717,0.3992983) (2.5432942,0.3170940) (2.5734456,0.2343946) (2.5995790,0.1513291) (2.6216533,0.0680274) (2.6396343,-0.0153805) (2.6534939,-0.0987644) (2.6632103,-0.1819941) (2.6687684,-0.2649396) (2.6701595,-0.3474715) (2.6673815,-0.4294608) (2.6604387,-0.5107796) (2.6493420,-0.5913009) (2.6341086,-0.6708989) (2.6147624,-0.7494493) (2.5913335,-0.8268297) (2.5638586,-0.9029190) (2.5323805,-0.9775985) (2.4969484,-1.0507516) (2.4576176,-1.1222640) (2.4144495,-1.1920242) (2.3675116,-1.2599232) (2.3168770,-1.3258550) (2.2626250,-1.3897166) (2.2048401,-1.4514083) (2.1436126,-1.5108339) (2.0790381,-1.5679004) (2.0112175,-1.6225189) (1.9402566,-1.6746039) (1.8662662,-1.7240743) (1.7893618,-1.7708528) (1.7096636,-1.8148662) (1.6272959,-1.8560460) (1.5423874,-1.8943278) (1.4550706,-1.9296518) (1.3654820,-1.9619629) (1.2737613,-1.9912106) (1.1800517,-2.0173493) (1.0844997,-2.0403381) (0.9872543,-2.0601412) (0.8884675,-2.0767276) (0.7882934,-2.0900714) (0.6868885,-2.1001519) (0.5844110,-2.1069532) (0.4810212,-2.1104647) (0.3768802,-2.1106811) (0.2721508,-2.1076018) (0.1669965,-2.1012318) (0.0615814,-2.0915809) (-0.0439299,-2.0786643) (-0.1493725,-2.0625021) (-0.2545819,-2.0431195) (-0.3593938,-2.0205469) (-0.4636446,-1.9948194) (-0.5671714,-1.9659773) (-0.6698127,-1.9340655) (-0.7714081,-1.8991338) (-0.8717991,-1.8612370) (-0.9708289,-1.8204340) (-1.0683428,-1.7767886) (-1.1641886,-1.7303690) (-1.2582167,-1.6812476) (-1.3502803,-1.6295011) (-1.4402355,-1.5752104) (-1.5279420,-1.5184602) (-1.6132628,-1.4593390) (-1.6960646,-1.3979393) (-1.7762183,-1.3343569) (-1.8535986,-1.2686910) (-1.9280847,-1.2010443) (-1.9995604,-1.1315222) (-2.0679139,-1.0602334) (-2.1330387,-0.9872891) (-2.1948330,-0.9128034) (-2.2532003,-0.8368924) (-2.3080495,-0.7596747) (-2.3592950,-0.6812708) (-2.4068567,-0.6018032) (-2.4506604,-0.5213960) (-2.4906377,-0.4401747) (-2.5267261,-0.3582661) (-2.5588694,-0.2757981) (-2.5870173,-0.1928995) (-2.6111258,-0.1096997) (-2.6311574,-0.0263286) (-2.6470808,0.0570836) (-2.6588710,0.1404067) (-2.6665098,0.2235105) (-2.6699851,0.3062653) (-2.6692915,0.3885420) (-2.6644302,0.4702120) (-2.6554087,0.5511478) (-2.6422410,0.6312231) (-2.6249479,0.7103128) (-2.6035561,0.7882934)};\draw[curve]plot coordinates{(-1.3017781,0.3941467) (-1.2890496,0.4325216) (-1.2743085,0.4702211) (-1.2575777,0.5071865) (-1.2388833,0.5433600) (-1.2182546,0.5786850) (-1.1957237,0.6131065) (-1.1713258,0.6465708) (-1.1450990,0.6790254) (-1.1170843,0.7104199) (-1.0873254,0.7407051) (-1.0558688,0.7698338) (-1.0227635,0.7977605) (-0.9880614,0.8244416) (-0.9518165,0.8498354) (-0.9140854,0.8739024) (-0.8749271,0.8966048) (-0.8344027,0.9179072) (-0.7925755,0.9377765) (-0.7495108,0.9561815) (-0.7052758,0.9730936) (-0.6599396,0.9884863) (-0.6135730,1.0023356) (-0.5662484,1.0146199) (-0.5180396,1.0253199) (-0.4690220,1.0344191) (-0.4192720,1.0419031) (-0.3688675,1.0477603) (-0.3178869,1.0519816) (-0.2664101,1.0545603) (-0.2145172,1.0554925) (-0.1622894,1.0547766) (-0.1098082,1.0524138) (-0.0571556,1.0484078) (-0.0044137,1.0427648) (0.0483350,1.0354937) (0.1010083,1.0266058) (0.1535239,1.0161149) (0.2057998,1.0040376) (0.2577543,0.9903925) (0.3093064,0.9752010) (0.3603756,0.9584869) (0.4108820,0.9402762) (0.4607469,0.9205974) (0.5098924,0.8994812) (0.5582418,0.8769606) (0.6057196,0.8530707) (0.6522516,0.8278488) (0.6977651,0.8013343) (0.7421892,0.7735686) (0.7854545,0.7445951) (0.8274934,0.7144590) (0.8682402,0.6832074) (0.9076314,0.6508890) (0.9456054,0.6175544) (0.9821030,0.5832555) (1.0170671,0.5480459) (1.0504432,0.5119806) (1.0821792,0.4751159) (1.1122254,0.4375093) (1.1405351,0.3992197) (1.1670639,0.3603067) (1.1917706,0.3208311) (1.2146164,0.2808546) (1.2355657,0.2404396) (1.2545858,0.1996492) (1.2716471,0.1585470) (1.2867228,0.1171973) (1.2997895,0.0756646) (1.3108267,0.0340137) (1.3198172,-0.0076902) (1.3267469,-0.0493822) (1.3316051,-0.0909970) (1.3343842,-0.1324698) (1.3350798,-0.1737357) (1.3336908,-0.2147304) (1.3302194,-0.2553898) (1.3246710,-0.2956504) (1.3170543,-0.3354494) (1.3073812,-0.3747247) (1.2956667,-0.4134148) (1.2819293,-0.4514595) (1.2661902,-0.4887992) (1.2484742,-0.5253758) (1.2288088,-0.5611320) (1.2072248,-0.5960121) (1.1837558,-0.6299616) (1.1584385,-0.6629275) (1.1313125,-0.6948583) (1.1024200,-0.7257042) (1.0718063,-0.7554169) (1.0395191,-0.7839502) (1.0056087,-0.8112594) (0.9701283,-0.8373020) (0.9331331,-0.8620372) (0.8946809,-0.8854264) (0.8548318,-0.9074331) (0.8136479,-0.9280230) (0.7711937,-0.9471639) (0.7275353,-0.9648259) (0.6827410,-0.9809814) (0.6368806,-0.9956053) (0.5900259,-1.0086746) (0.5422498,-1.0201690) (0.4936271,-1.0300706) (0.4442337,-1.0383638) (0.3941467,-1.0450357) (0.3434442,-1.0500759) (0.2922055,-1.0534766) (0.2405106,-1.0552324) (0.1884401,-1.0553405) (0.1360754,-1.0538009) (0.0834982,-1.0506159) (0.0307907,-1.0457905) (-0.0219649,-1.0393321) (-0.0746862,-1.0312510) (-0.1272910,-1.0215598) (-0.1796969,-1.0102734) (-0.2318223,-0.9974097) (-0.2835857,-0.9829886) (-0.3349063,-0.9670327) (-0.3857041,-0.9495669) (-0.4358996,-0.9306185) (-0.4854144,-0.9102170) (-0.5341714,-0.8883943) (-0.5820943,-0.8651845) (-0.6291084,-0.8406238) (-0.6751401,-0.8147506) (-0.7201178,-0.7876052) (-0.7639710,-0.7592301) (-0.8066314,-0.7296695) (-0.8480323,-0.6989697) (-0.8881091,-0.6671785) (-0.9267993,-0.6343455) (-0.9640424,-0.6005221) (-0.9997802,-0.5657611) (-1.0339570,-0.5301167) (-1.0665194,-0.4936446) (-1.0974165,-0.4564017) (-1.1266002,-0.4184462) (-1.1540248,-0.3798373) (-1.1796475,-0.3406354) (-1.2034284,-0.3009016) (-1.2253302,-0.2606980) (-1.2453188,-0.2200874) (-1.2633631,-0.1791331) (-1.2794347,-0.1378991) (-1.2935086,-0.0964498) (-1.3055629,-0.0548498) (-1.3155787,-0.0131643) (-1.3235404,0.0285418) (-1.3294355,0.0702033) (-1.3332549,0.1117552) (-1.3349925,0.1531327) (-1.3346458,0.1942710) (-1.3322151,0.2351060) (-1.3277043,0.2755739) (-1.3211205,0.3156115) (-1.3124739,0.3551564) (-1.3017781,0.3941467)};\draw[thick](-.5,2.74)--(3,.64);\fill(0,0)circle(1.5pt)node[below left]{$(3,6)$};\fill(1.3880239520958084,1.6071856287425148)circle(1.7pt)node[above right]{$P_*$};\draw[axis](1.3880239520958084,1.6071856287425148)--(0.9380239520958085,0.8571856287425148)node[below left]{$\nabla P$};
\end{tikzpicture}
\caption{最外层所画等值线与约束直线相切；该处梯度指向无约束高值区。}\label{fig:c3-e5-s}
\end{figure}
\end{solution}

\begin{example}[单纯形上的熵极值]\label{ex:c3-e6}
在$x,y,z\ge0,x+y+z=1$上求$H=-x\ln x-y\ln y-z\ln z$的最大最小值，约定$0\ln0=0$。推广到$n$个分量。 见图~\ref{fig:c3-e6-q}。
\end{example}

\begin{figure}[H]\centering
\begin{tikzpicture}[>=Stealth,line cap=round,line join=round]
\draw[curve](0,0)--(4,0)--(1.4,2.7)--cycle;\node[below]at(0,0){$(1,0,0)$};\node[below]at(4,0){$(0,1,0)$};\node[above]at(1.4,2.7){$(0,0,1)$};\fill(1.8,.9)circle(1.5pt)node[right]{$C$};
\end{tikzpicture}
\caption{可行域是闭三角形，边界也属于定义域。}\label{fig:c3-e6-q}
\end{figure}
\begin{solution}
允许分量为零时，约束集合是包含边和顶点的闭单纯形。图~\ref{fig:c3-e6-q}中的三角形内部适用普通乘子法，边界处对数导数发散，但$-x\ln x$本身连续趋零，所以边界仍是必须比较的可行状态。


对 $\phi(u)=-u\ln u$ 用乘积法则，$\phi'=-(\ln u+u/u)=-\ln u-1$，再求导得 $\phi''=-1/u<0$。内点拉格朗日方程为 $-\ln x-1=\lambda$，对 $y,z$ 同理，因此 $\ln x=\ln y=\ln z$；再用总和为一，得 $x=y=z=1/3$、$H=\ln3$。边界的两个分量情形在下面单独比较。

推广到$n$个分量时，可直接用上述严格凹性证明全局上界。记$\bar x=1/n$，由于$\phi'$严格递减，对任意$u>0$，
\[
 \phi(u)-\phi(\bar x)
 =\int_{\bar x}^{u}\phi'(s)\dd s
 \le\phi'(\bar x)(u-\bar x).
\]
当$u<\bar x$时，积分方向与$\phi'(s)-\phi'(\bar x)$的符号同时改变，不等号方向仍如上；当$u=0$时由连续性取极限。除$u=\bar x$外不等式严格。依次令$u=x_i$后相加，线性项因$\sum_i(x_i-\bar x)=0$相消，所以
\[\frac1n\sum_{i=1}^n\phi(x_i)\le\phi\!\left(\frac1n\sum_i x_i\right)=\frac{\ln n}{n}.
\]
故$H\le\ln n$，等号仅在所有$x_i=1/n$时取得。另一方面$0\le x_i\le1$，每项$-x_i\ln x_i\ge0$，只有一个分量为$1$的顶点可使总和为零。因此
\[\boxed{H_{\max}=\ln n,\quad H_{\min}=0}.\]
图~\ref{fig:c3-e6-s}沿三角形中线取$x=y=(1-z)/2$，在$z=1/3$达到最高值。该截线显示等分点附近的变化；整个单纯形上的极值已由上述凹性与边界分析确定。


三分量边界上的比较可直接完成。例如$z=0$时，令$y=1-x$，
\[
\begin{aligned}
H_b(x)&=-x\ln x-(1-x)\ln(1-x),\\
H_b'&=-(\ln x+1)+\ln(1-x)+1=\ln\frac{1-x}{x},\\
H_b''&=-\frac1{1-x}-\frac1x=-\frac1{x(1-x)}<0.
\end{aligned}
\]
所以边内最高值为$H_b(1/2)=\ln2$，两端值均为零。三条边相同，而$\ln2<\ln3$，故内点等分解超过全部边界。解答图沿$x=y$的截线还包含一个二分量边界点$z=0$，该点高度为$\ln2$，不是零。

若只剩 $k<n$ 个正分量，即使在它们内部等分，熵也只有 $\ln k<\ln n$，仍不能达到全单纯形的最大值。

最小值的等号条件须逐项成立：总和为零要求每个$x_i$为$0$或$1$，总和约束则只允许一个$1$。边界点若仍有多个正分量，熵便大于零。

\begin{figure}[H]\centering
\begin{tikzpicture}[>=Stealth,line cap=round,line join=round]
\begin{scope}[x=2.25080cm,y=1.28571cm]\draw[axis](0,0)--(1.2469999999999999,0)node[right]{$z$};\draw[axis](0,0)--(0,1.4)node[above]{$H$};\draw[curve,domain=0.003:0.997,samples=101]plot(\x,{-(1-\x)*ln((1-\x)/2)-\x*ln(\x)});\draw[guide](.333333,0)--(.333333,1.098612);\node[below]at(.333333,0){$1/3$};\end{scope}
\end{tikzpicture}
\caption{在$x=y$的截面中，熵在三个分量相等时最大。}\label{fig:c3-e6-s}
\end{figure}
\end{solution}


\section{梯度、散度与旋度}\label{sec:c3-fields}
局部场的方向变化率、通量密度和环量密度分别对应不同算子。各算子从几何定义建立，再用于势能、连续性及环路方程。

\subsection{方向导数与正交坐标}
沿单位向量$\vct n$的方向导数为$D_{\vct n}f=\nabla f\cdot\vct n$。若方向由非单位向量$\vct a$给出，先除以$|\vct a|$；否则算得的是参数变化率，并非单位长度的变化率。梯度指向增长最快的方向，最大方向导数为$|\nabla f|$。等值面上$\dd f=0$，故梯度垂直于等值面。

正交曲线坐标$q_i$的尺度因子为$h_i=|\partial\vct r/\partial q_i|$，线元$\dd\vct r=\sum_i h_i\vct e_i\dd q_i$。由$\dd f=\nabla f\cdot\dd\vct r$比较系数，得
\[\nabla f=\sum_i\vct e_i\frac1{h_i}\frac{\partial f}{\partial q_i}.\]
柱坐标$(r,\varphi,z)$满足$x=r\cos\varphi,y=r\sin\varphi$，尺度因子$(1,r,1)$。球坐标$(r,\theta,\varphi)$取$\theta$为从正$z$轴量起的极角，尺度因子$(1,r,r\sin\theta)$。两种坐标中的$r$含义不同，题内需保持一致。

\subsection{方向导数的长度尺度与曲线坐标}
沿经过$P$的单位方向$\vct n$走长度$s$，路径为$\vct r=P+s\vct n$。可微性给出$f(P+s\vct n)-f(P)=s\nabla f(P)\cdot\vct n+o(s)$，所以方向导数是梯度的投影。Cauchy不等式给出$|D_{\vct n}f|\le|\nabla f|$，沿梯度方向取最大增长率，反向取最大下降率。

若用任意参数$u$描述路径，则$\dd f/\dd u=\nabla f\cdot\vct r_u$。将$\vct r_u$写成速率$|\vct r_u|$乘单位切向，便得到参数变化率与单位长度变化率之间的转换。习题~\ref{ex:c3-h1}中的角速度需乘半径才成为速率；只把$\omega$乘方向导数，量纲就不一致。

曲线坐标的尺度因子把坐标增量换算成实际长度。固定其他坐标改变$q_i$时，实际小位移长度为$h_i\dd q_i$，所以梯度沿该单位基矢的分量为$f_{q_i}/h_i$。在柱坐标中，角度改变$\dd\varphi$扫过弧长$r\dd\varphi$；在球坐标中，方位角扫过的圆半径是$r\sin\theta$，所以弧长为$r\sin\theta\dd\varphi$。公式中的几何因子由这些长度直接确定。

坐标基矢本身也随位置变化。柱坐标中$\partial_\varphi\vct e_r=\vct e_\varphi$、$\partial_\varphi\vct e_\varphi=-\vct e_r$。向量场求导时，既要对分量求导，也要计入基矢转动。把柱坐标分量直接代进直角坐标公式，会丢掉这部分几何变化。

\subsection{由通量与环量定义算子}
对小闭合体积的外向通量除以体积，再取极限，得到散度：
\[\nabla\cdot\vct A=\lim_{\Delta V\to0}\frac1{\Delta V}\iint_{\partial\Delta V}\vct A\cdot\vct n\dd S.
\]
小体内若有净流出，散度为正。旋度沿单位法向$\vct n$的分量，则由与该法向按右手规则定向的小回路环量给出：
\[(\nabla\times\vct A)\cdot\vct n
=\lim_{\Delta S\to0}\frac1{\Delta S}\oint_{\partial\Delta S}\vct A\cdot\dd\vct l.
\]
同一向量场可以散度为零而旋度非零。定义中的边界方向与法向必须相配，换向会改变环量的符号。

散度的直角坐标式可由小长方体求出。左右两面的净$x$向通量为$[A_x(x+\Delta x)-A_x(x)]\Delta y\Delta z\simeq(\partial_xA_x)\Delta V$，另外两对面同理。三项相加并除以体积，就得到$\partial_xA_x+\partial_yA_y+\partial_zA_z$。正值表示局部净流出，负值表示净流入。

旋度则从小矩形环路求得。以$+z$为法向，边界从上方看逆时针。上下两条$x$向边贡献$-\partial_yA_x\Delta x\Delta y$，左右两条$y$向边贡献$\partial_xA_y\Delta x\Delta y$，相加后除以面积得到$(\nabla\times\vct A)_z=\partial_xA_y-\partial_yA_x$。交换法向或回路方向会反转符号，保持二者按右手规则配对则不改变物理结论。

梯度作用于标量，散度把向量变成标量，旋度把向量变成向量。其量纲均比原场多一个“每单位长度”，但含义不同：梯度描述标量变化，散度描述通量收支，旋度描述有向环量密度。刚体转动场可以不可压缩却有旋，径向点源场可以在源外无散度却有包围通量。

图~\ref{fig:c3-flux-difference}只画垂直于$x$轴的两面。两面的外法向相反，故净流出量是右面值减左面值，首项为$\partial_xA_x\,\dd V$。另外两对面按相同规则相加，就得到散度。
\begin{figure}[H]\centering
\begin{tikzpicture}[>=Stealth,line cap=round,line join=round]
\fill[gray!18](0,0)--(.7,.45)--(.7,1.95)--(0,1.5)--cycle;
\fill[gray!18](3.2,0)--(3.9,.45)--(3.9,1.95)--(3.2,1.5)--cycle;
\draw[curve](0,0)--(.7,.45)--(.7,1.95)--(0,1.5)--cycle;
\draw[curve](3.2,0)--(3.9,.45)--(3.9,1.95)--(3.2,1.5)--cycle;
\draw[curve](0,0)--(3.2,0) (0,1.5)--(3.2,1.5) (.7,1.95)--(3.9,1.95);\draw[guide](.7,.45)--(3.9,.45);
\draw[->,thick](.35,.975)--(-1,.975)node[left]{$-\vct e_x$};\draw[->,thick](3.55,.975)--(4.9,.975)node[right]{$+\vct e_x$};
\draw[<->](0,-.55)--(3.2,-.55)node[midway,below]{$\dd x$};\node[above]at(.35,2.2){左面：$x$};\node[above]at(3.55,2.2){右面：$x+\dd x$};
\node[above]at(1.9,2.8){两阴影面的面积均为$\dd y\dd z$};
\node[below]at(1.6,-1.25){$\dd\Phi_x=[A_x(x+\dd x)-A_x(x)]\dd y\dd z$};
\end{tikzpicture}
\caption{散度来自各对相反外法向面上的通量差。}\label{fig:c3-flux-difference}
\end{figure}

\subsection{直角坐标公式}
\[\nabla f=(f_x,f_y,f_z),\qquad \nabla\cdot\vct A=A_{x,x}+A_{y,y}+A_{z,z},\]
\[\nabla\times\vct A=(A_{z,y}-A_{y,z},\ A_{x,z}-A_{z,x},\ A_{y,x}-A_{x,y}),\]
\[\nabla^2f=\nabla\cdot\nabla f=f_{xx}+f_{yy}+f_{zz}.\]
连续二阶偏导保证$\nabla\times\nabla f=0$和$\nabla\cdot(\nabla\times\vct A)=0$。这些恒等式描述局部关系。场域有孔洞或奇点时，还须另查绕孔闭合回路的环量。

\subsection{柱坐标公式}
这里$A_r,A_\varphi,A_z$均为单位基矢的物理分量。
\[\nabla f=\vct e_r f_r+\vct e_\varphi f_\varphi/r+\vct e_z f_z,\]
\[\nabla\cdot\vct A=\frac1r\partial_r(rA_r)+\frac1r\partial_\varphi A_\varphi+\partial_z A_z,\]
\begin{align*}
(\nabla\times\vct A)_r&=\frac1r\partial_\varphi A_z-\partial_z A_\varphi,\\
(\nabla\times\vct A)_\varphi&=\partial_z A_r-\partial_r A_z,\\
(\nabla\times\vct A)_z&=\frac1r[\partial_r(rA_\varphi)-\partial_\varphi A_r],\\
\nabla^2f&=\frac1r\partial_r(rf_r)+\frac1{r^2}f_{\varphi\varphi}+f_{zz}.
\end{align*}
例如$\vct e_r$随$\varphi$改变方向，把柱坐标分量逐项当成直角坐标分量求导会遗漏尺度和基矢变化。上式中的$r$因子正是由这些几何变化产生的。

\subsection{球坐标公式}
\[\nabla f=\vct e_r f_r+\vct e_\theta f_\theta/r+\vct e_\varphi f_\varphi/(r\sin\theta),\]
\[\nabla\cdot\vct A=\frac1{r^2}\partial_r(r^2A_r)+\frac1{r\sin\theta}\partial_\theta(\sin\theta A_\theta)+\frac1{r\sin\theta}\partial_\varphi A_\varphi.
\]
\begin{align*}
(\nabla\times\vct A)_r&=\frac{\partial_\theta(\sin\theta A_\varphi)-\partial_\varphi A_\theta}{r\sin\theta},\\
(\nabla\times\vct A)_\theta&=\frac{\partial_\varphi A_r/\sin\theta-\partial_r(rA_\varphi)}r,\\
(\nabla\times\vct A)_\varphi&=\frac{\partial_r(rA_\theta)-\partial_\theta A_r}r,\\
\nabla^2f&=\frac1{r^2}\partial_r(r^2f_r)+\frac1{r^2\sin\theta}\partial_\theta(\sin\theta f_\theta)+\frac{f_{\varphi\varphi}}{r^2\sin^2\theta}.
\end{align*}
在轴线或极点处出现的分母为零，可能只是坐标退化。应回到直角坐标或取物理场的连续极限，而不应把每个坐标奇点都解释成真实源。

拉普拉斯算子是先取梯度、再取散度。对球对称势$f(r)$，梯度沿径向，球面面积随$r^2$变化，所以无源条件要求$r^2f'(r)$不随半径变化，积分给出$A/r+B$。柱对称势的单位长度柱面面积随$r$变化，因而要求$rf'(r)$不变，得到$A\ln r+B$。习题~\ref{ex:c3-h11}将这两种几何守恒形式直接积分。

坐标轴、极点和真实场奇点须区分。球坐标在$\theta=0,\pi$处方位角失去唯一性，分量公式中的$1/\sin\theta$可能表现为坐标退化；一个在直角坐标中光滑的场仍应有有限且一致的物理极限。点电荷原点或涡旋轴线则是真实的场域排除点，不能以换坐标消除。使用积分定理时，后者必须明确挖去或以源项补充。

\subsection{电磁场中的局部方程}
在宏观SI制中，以下$\rho,\vct J$分别表示自由电荷密度与自由电流密度：
\[\nabla\cdot\vct D=\rho,\quad \nabla\cdot\vct B=0,\quad
\nabla\times\vct E=-\partial_t\vct B,\quad
\nabla\times\vct H=\vct J+\partial_t\vct D.
\]
真空中$\vct D=\varepsilon_0\vct E,\vct B=\mu_0\vct H$，所以最后一式也写为$\nabla\times\vct B=\mu_0\vct J+\mu_0\varepsilon_0\partial_t\vct E$。将两种写法混用会遗漏$\mu_0$。对最后一式取散度，并利用$\nabla\cdot\vct D=\rho$，得到连续性方程$\partial_t\rho+\nabla\cdot\vct J=0$。下一讲用面积分和体积分把这些局部关系还原为有限区域的通量与环量关系。

\Needspace{8\baselineskip}
宏观场方程中$\vct D,\vct H$把介质束缚源的一部分吸收到场的定义中，所以右侧采用自由电荷与自由电流。真空中没有这层区分，才可直接换成$\vct E,\vct B$及$\varepsilon_0,\mu_0$。不同记号的公式不能只替换场符号而保留原系数。

位移电流项的系数可由电荷连续性核对。对$\nabla\times\vct H=\vct J+\partial_t\vct D$取散度，左侧恒为零，右侧成为$\nabla\cdot\vct J+\partial_t\rho$。若删去时间项，一般随时间变化的自由电荷便无法满足连续性。使用静态近似时，应说明相应时间变化为何可以忽略。

静电势满足$\vct E=-\nabla\Phi$，因此场线沿电势下降方向。梯度求场时应先说明所用坐标及角度定义；球坐标中$E_\theta=-(1/r)\partial_\theta\Phi$，方位分量为$-(r\sin\theta)^{-1}\partial_\varphi\Phi$。例~\ref{ex:c3-e10}用同一个偶极势推导分量和场线，角度方向在每一步保持一致。

\begin{example}[点电荷、偶极势与场线]\label{ex:c3-e10}
正负点电荷$\pm q$相距$\ell$，偶极矩$\vct p=q\ell\vct e_z$由负电荷指向正电荷。在$r\gg\ell$处求电势、电场及场线方程；先说明点电荷势的梯度。 见图~\ref{fig:c3-e10-q}。
\end{example}

\begin{figure}[H]\centering
\begin{tikzpicture}[>=Stealth,line cap=round,line join=round]
\draw[axis](0,-1.7)--(0,2.15)node[above]{$z$};
\draw[axis](-.3,0)--(3.05,0)node[right]{$x$};
\fill(0,.75)circle(2pt)node[above right]{$+q$};\fill(0,-.75)circle(2pt)node[right]{$-q$};
\draw[axis](-.65,-.75)--(-.65,.75)node[midway,left]{$\vct p$};
\draw[axis](0,0)--(2.65,1.65)node[above]{$P$};
\node[above,sloped]at(1.85,1.15){$\vct r$};
\draw(0,.48)arc(90:31.9:.48);\node at(.26,.57){$\theta$};
\node[below left]at(0,0){$O$};\draw[<->](.7,-.75)--(.7,.75)node[midway,right]{$\ell$};
\end{tikzpicture}
\caption{两电荷置于同一$z$轴上。极角$\theta$在正$z$轴与位置矢量之间量取，偶极矩由负电荷指向正电荷。}\label{fig:c3-e10-q}
\end{figure}
\begin{solution}
先由点电荷势的径向导数求场，再对正负电荷的势作远场展开。两电荷的单极项抵消，保留首个非零偶极项。图~\ref{fig:c3-e10-q}固定了极角从正$z$轴量起以及偶极矩方向，电场分量与场线积分都要按这一约定进行。


记$k=1/(4\pi\varepsilon_0)$。点电荷势为$\Phi=kQr^{-1}$，径向导数是$\partial_r\Phi=-kQr^{-2}$，故$\vct E=-\nabla\Phi=kQ\vct e_r/r^2$。对偏离原点的电荷，令$a=|\vct a|$，先提出观测距离$r$：
\[
 \frac1{|\vct r-\vct a|}=\frac1r(1+\zeta)^{-1/2},
 \qquad \zeta=-\frac{2\vct a\cdot\vct e_r}{r}+\frac{a^2}{r^2}.
\]
当$a/r\ll1$，用$(1+\zeta)^{-1/2}=1-\zeta/2+O(\zeta^2)$，保留$a/r$的一阶项，得到
\[\frac1{|\vct r-\vct a|}=\frac1r+\frac{\vct a\cdot\vct e_r}{r^2}+O(a^2/r^3).
\]
取正电荷位于 $\vct a=\ell\vct e_z/2$，负电荷在 $-\vct a$，两项的单极部分抵消，偶极部分相加为 $kq(\ell\cos\theta)/r^2=kp\cos\theta/r^2$。由此求电场：
\[
E_r=-\partial_r(kp\cos\theta\,r^{-2})=2kp\cos\theta/r^3,
\]
\[
E_\theta=-\frac1r\partial_\theta(kp\cos\theta/r^2)
=-\frac1r\frac{kp}{r^2}(-\sin\theta)=kp\sin\theta/r^3.
\]
电势与场分量为
\[\boxed{\Phi=\frac{kp\cos\theta}{r^2}},\qquad
\boxed{E_r=\frac{2kp\cos\theta}{r^3},\quad E_\theta=\frac{kp\sin\theta}{r^3},\quad E_\varphi=0}.
\]
球坐标梯度的角分量含$1/r$，若漏掉它会使两个电场分量量纲不一致。图~\ref{fig:c3-e10-s}中的场线切向与$\vct E$平行，因$\dd\vct r=\vct e_r\dd r+\vct e_\theta r\dd\theta$，有
\[\frac{\dd r}{r\dd\theta}=\frac{E_r}{E_\theta}=2\cot\theta.
\]
在 $0<\theta<\pi$ 内选线上的一点 $(r_*,\theta_*)$，分离变量积分。右侧令 $u=\sin\theta$、$\dd u=\cos\theta\dd\theta$，得到
\[
\int_{r_*}^r\frac{\dd s}s=2\int_{\theta_*}^\theta\frac{\cos v}{\sin v}\dd v
=2[\ln\sin v]_{\theta_*}^\theta,
\qquad \ln\frac r{r_*}=2\ln\frac{\sin\theta}{\sin\theta_*}.
\]
取指数，并记 $C=r_*/\sin^2\theta_*$，得到 $\boxed{r=C\sin^2\theta}$。
$C>0$标识不同场线，方位角沿线保持不变。该方程属于理想点偶极远场，不能一直用到实际两点电荷之间。轴线上$\sin\theta=0$时，除法推导不适用，轴线场应单独由电场方向解释。


题图中的极角必须从正$z$轴量起。由$\cos\theta=z/r$，偶极势也可写成$\Phi=kpz/r^3$；在赤道面$z=0$，$\Phi=0$而电场并不为零，此处$\vct E=-kp\vct e_z/r^3$，与解答图中场线向下穿过赤道一致。轴线上则$\vct E$均沿正$z$方向：正半轴向外、负半轴向内。

对于有限间距的两电荷，以上等式表示偶极近似的首项。两电荷关于原点反对称，二次阶势项相消，下一个非零势项为$O(kq\ell^3/r^4)$。因此场线在离电荷距离与$\ell$相当的区域必须改用完整两电荷场，图中原点附近的曲线仅表示理想偶极的极限形状。

对 $0<\theta<\pi$，$E_\theta$ 为正，表示场线沿极角增加方向从上半空间穿向下半空间。场线微小位移的径向和极角向长度分别为 $\dd r$ 与 $r\dd\theta$，两者比值才与场分量比值相等。

图~\ref{fig:c3-e10-s}中的$C$是每条场线的赤道最大半径：在$\theta=\pi/2$时$r=C$，而两侧$r$逐渐减小。

\begin{figure}[H]\centering
\begin{tikzpicture}[>=Stealth,line cap=round,line join=round]
\draw[axis](0,-2)--(0,2)node[above]{$z$};\foreach\c in{1,1.5,2}{\draw[curve,domain=2:178,samples=121]plot({\c*sin(\x)^3},{\c*sin(\x)^2*cos(\x)});\draw[curve,domain=2:178,samples=121]plot({-\c*sin(\x)^3},{\c*sin(\x)^2*cos(\x)});}\fill(0,0)circle(2pt);\draw[axis](1.5,.12)--(1.5,-.12);\draw[axis](-1.5,.12)--(-1.5,-.12);\node[right]at(2.1,.6){$r=C\sin^2\theta$};
\end{tikzpicture}
\caption{子午面场线在赤道处指向负$z$方向，与$E_\theta>0$一致。}\label{fig:c3-e10-s}
\end{figure}
\end{solution}

\section{二次变分与曲线坐标的几何}
稳定性与几何尺度都会影响二阶量。先把允许扰动写在同一组独立坐标中，再用实际长度、面积和体积构造二次变化，坐标分量不能直接当成物理长度。

\subsection{约束极值的二阶判别}
设$g(\vct r)=0$，在正规驻点$\nabla f=\lambda\nabla g$附近，取可行曲线$\vct r(s)$，令$\vct t=\vct r'(0)$、$\vct a=\vct r''(0)$。一次求导给出$\nabla g\cdot\vct t=0$；二次用链式和乘积法则，得到
\[0=\left.\frac{\dd^2g}{\dd s^2}\right|_0
=\sum_{ij}g_{ij}t_it_j+\sum_i g_i a_i
=\vct t^{\mathsf T}H_g\vct t+\nabla g\cdot\vct a.\]
目标函数的二阶导数也有两部分：
\[\left.\frac{\dd^2f}{\dd s^2}\right|_0
=\vct t^{\mathsf T}H_f\vct t+\nabla f\cdot\vct a.\]
用乘子关系将后一项改为$\lambda\nabla g\cdot\vct a=-\lambda\vct t^{\mathsf T}H_g\vct t$，得到
\[\left.\frac{\dd^2 f}{\dd s^2}\right|_0
=\vct t^{\mathsf T}(H_f-\lambda H_g)\vct t.
\]
因此需要在约束切空间上考察拉格朗日函数的Hessian，而不是在整个空间考察$H_f$。约束的曲率会通过$-\lambda H_g$进入二阶项。多个约束时将各项相加；可行曲线只在交线切向上变化，判别空间随之缩小。

可行曲线的局部展开是$\vct r(s)=\vct r_0+s\vct t+s^2\vct a/2+\cdots$。一阶位移沿切线，二阶位移一般还含为保持约束所需的法向分量；上面的$-\lambda H_g$正是这部分曲率的贡献。只看$H_f$会漏掉它，而要求修正矩阵在整个空间正定又过强，实际只需在切空间判别。球面重力势在直角坐标中Hessian为零，却仍在南极有稳定平衡，习题~\ref{ex:c3-h14}给出了这一应用。

若目标是势能，小振动还需同样投影动能到切向坐标。两个切向方向的二阶势能系数与相应质量共同决定频率。出现零系数时，须分析高阶束缚或自由漂移，不能仅由“势能不下降”就断言完整位置--速度状态稳定。

\subsection{正交坐标中的散度微元}
在正交坐标$q_i$中，一个小体的体积为$h_1h_2h_3\dd q_1\dd q_2\dd q_3$。与第一方向垂直的两面面积为$h_2h_3\dd q_2\dd q_3$。两面通量之差的一阶项是$\partial_{q_1}(h_2h_3A_1)\dd q_1\dd q_2\dd q_3$，三方向相加再除以体积，
\[\nabla\cdot\vct A=\frac1{h_1h_2h_3}\sum_i\partial_{q_i}\left(\frac{h_1h_2h_3}{h_i}A_i\right).
\]
代入柱坐标尺度$(1,r,1)$和球坐标尺度$(1,r,r\sin\theta)$，就分别得到前面的公式。多出的几何因子来自相邻坐标面的面积变化。对旋度沿第$i$方向的分量，用另外两个坐标方向张成的小回路同样计算，得到循环指标形式
\[(\nabla\times\vct A)_i=\frac{\partial_{q_j}(h_kA_k)-\partial_{q_k}(h_jA_j)}{h_jh_k}.
\]
这里$(i,j,k)$按右手正交次序循环。梯度、散度和旋度均可从长度、面积、体积的几何定义推导，无需把三套公式作为互不相关的记忆项。

\subsection{守恒场的局部条件与整体条件}
若$\vct F=-\nabla U$且$U$单值，沿光滑路径$\vct r(s)$有
\[\vct F\cdot\frac{\dd\vct r}{\dd s}
=-\nabla U\cdot\frac{\dd\vct r}{\dd s}=-\frac{\dd U}{\dd s}.\]
从起点$A$积到终点$B$，功为$-[U]_A^B=U(A)-U(B)$；闭路的两端相同，所以功为零。

静电场、无旋稳恒理想流与中心力问题常能使用势函数，但场的奇点必须从定义域中明确分离。例如点电荷场在原点以外无散度，仍有非零的包围通量；涡旋场在轴外无旋，仍有非零的绕轴环量。局部微分式与有限边界积分反映不同层次的信息，两者通过带正则性条件的积分定理相接。

局部无旋与全域存在单值势之间还相隔一个路径条件。若所有闭合回路环量都为零，固定参考点后可用路径积分定义势；两条到达同一点的路径之差是一条闭路，所以势与选路无关。在单连通且场光滑的区域，无旋条件可保证这一性质；绕孔区域中的闭路却未必能连续收缩到一点，须另查绕孔环量。

同样，点源外散度为零只说明任一不包含源的小体积没有净源。围住原点的球面可以有固定非零通量，因为内部不再满足光滑无源条件。习题~\ref{ex:c3-h8}通过球壳内、外两个边界说明通量怎样抵消，习题~\ref{ex:c3-h10}则用绕轴场说明无旋条件为何只给出局部势。

应用积分定理时，除微分算子外还须给出区域边界、场的光滑范围和定向；这些条件决定局部微分关系能否用于整个区域。


\section{习题与解答}

\Needspace{8\baselineskip}
\begin{exercise}[圆周上的方向变化率]\label{ex:c3-h1}
温度场$T(x,y)=T_0+ax^2+by^2$。粒子沿半径$R$的圆以角速度$\omega$运动，求其测得的$\dot T$；再求点$(R/\sqrt2,R/\sqrt2)$处沿圆周正向的单位长度温度变化率。 见图~\ref{fig:c3-h1-q}。
\end{exercise}

\begin{figure}[H]\centering
\begin{tikzpicture}[>=Stealth,line cap=round,line join=round]
\draw[curve](0,0)circle(1.4);\draw[axis](-1.8,0)--(2,0)node[right]{$x$};\draw[axis](0,-1.7)--(0,2)node[above]{$y$};\fill(.98995,.98995)circle(1.5pt)node[right]{$P$};\draw[axis](.98995,.98995)--(.4,1.58);
\end{tikzpicture}
\caption{圆周正向取逆时针，单位切向与梯度不一定平行。}\label{fig:c3-h1-q}
\end{figure}
\begin{solution}
粒子沿等半径圆周运动，角速度给出位置参数，但测温变化率需要梯度与线速度的点积。图~\ref{fig:c3-h1-q}中的逆时针切向与径向垂直；在指定点先将切向归一化，才能求“每单位长度”的变化率。


取 $x=R\cos\omega t,y=R\sin\omega t$，求导为 $\dot x=-R\omega\sin\omega t$、$\dot y=R\omega\cos\omega t$。代入链式法则，
\[
\begin{aligned}
\dot T&=2ax\dot x+2by\dot y\\
&=-2aR^2\omega\cos\omega t\sin\omega t+2bR^2\omega\sin\omega t\cos\omega t\\
&=R^2\omega(b-a)\sin(2\omega t).
\end{aligned}
\]
也可先将轨迹代入温度，得到$T=T_0+aR^2\cos^2\omega t+bR^2\sin^2\omega t$再求导。图~\ref{fig:c3-h1-s}中$P$处的正向单位切向为$\vct n=(-1,1)/\sqrt2$，梯度为$\nabla T=\sqrt2R(a,b)$，因此
\[\boxed{D_{\vct n}T=R(b-a)}.
\]
该点沿圆周正向的有向速度为 $R\omega$，乘以正向单位长度变化率，得 $R^2\omega(b-a)$，与 $\omega t=\pi/4$ 时的链式法则一致；反向运动时 $\omega<0$，读数变化相应变号。若 $a=b$，整个圆上温度相同，梯度却未必为零，只是与圆周切向垂直。



\begin{figure}[H]\centering
\begin{tikzpicture}[>=Stealth,line cap=round,line join=round]
\draw[axis](0,0)--(3.2,0)node[right]{$x$};\draw[axis](0,0)--(0,3)node[above]{$y$};
\draw[guide](2,0)arc(0:90:2);\coordinate(P)at({sqrt(2)},{sqrt(2)});
\fill(P)circle(1.5pt)node[below right]{$P$};
\draw[axis](P)--({sqrt(2)-.75},{sqrt(2)+.75})node[left]{$\vct n$};
\draw[axis](P)--({sqrt(2)+.55},{sqrt(2)+1.1})node[right]{$\nabla T$};
\node[align=left,right]at(4,1.7){$P=(R/\sqrt2,R/\sqrt2)$\\[5pt]$D_{\vct n}T=R(b-a)$};
\end{tikzpicture}
\caption{图示取$b=2a>0$，圆弧、切向和梯度对应同一点$P$；一般参数的结果由点积公式给出。}\label{fig:c3-h1-s}
\end{figure}
\end{solution}

\Needspace{8\baselineskip}
\begin{exercise}[转动坐标中的链式法则]\label{ex:c3-h2}
令$x=s\cos t,y=s\sin t$，$w=f(x,y)$。求$w_s,w_t$，并证明在$s>0$处$w_s^2+w_t^2/s^2=f_x^2+f_y^2$。 见图~\ref{fig:c3-h2-q}。
\end{exercise}

\begin{figure}[H]\centering
\begin{tikzpicture}[>=Stealth,line cap=round,line join=round]
\draw[axis](0,0)--(3,0)node[right]{$x$};\draw[axis](0,0)--(0,2.5)node[above]{$y$};\draw[axis](0,0)--(2,1.4)node[midway,above]{$s$};\draw(1,0)arc(0:35:1);\node at(1.2,.4){$t$};
\end{tikzpicture}
\caption{变量$t$在此是角坐标，$s$是径向长度。}\label{fig:c3-h2-q}
\end{figure}
\begin{solution}
$s$是径向长度，$t$是角参数。固定角改变$s$与固定半径改变$t$对应不同长度尺度。图~\ref{fig:c3-h2-q}中的角弧长为$s\dd t$，因此两个方向导数在组成平方和前，需要把角参数导数除以$s$。


固定$t$时$(x_s,y_s)=(\cos t,\sin t)$；固定$s$时$(x_t,y_t)=s(-\sin t,\cos t)$。于是
\[w_s=f_x\cos t+f_y\sin t,\quad
w_t=s(-f_x\sin t+f_y\cos t).
\]
图~\ref{fig:c3-h2-s}中的两单位向量正交。直接将两式平方相加，第一式的交叉项为 $2f_xf_y\sin t\cos t$，第二式除以 $s^2$ 后的交叉项为其相反数；它们相消后，$f_x^2$、$f_y^2$ 的系数都是 $\sin^2t+\cos^2t=1$，故
\[\boxed{w_s^2+w_t^2/s^2=f_x^2+f_y^2}.\]
角方向的弧长微元为$s\dd t$，所以单位长度的变化率为$w_t/s$。在$s=0$处角坐标退化，商式不能直接使用；若$f$在原点可微，应通过直角坐标梯度讨论方向极限。


图中角向参数矢量 $(x_t,y_t)$ 除以 $s$ 后，便与径向矢量组成正交单位基。对应的梯度表达式为
\[\nabla f=\vct e_s w_s+\vct e_t\frac{w_t}{s}.\]
这里$t$在本题表示角坐标，并非时间，$w_t$本身不应解释为沿轨迹的时间导数。

正交变换可写成矩阵形式：
\[\binom{w_s}{w_t/s}=
\begin{pmatrix}\cos t&\sin t\\-\sin t&\cos t\end{pmatrix}
\binom{f_x}{f_y}.\]
该矩阵转置乘自身为单位阵，所以平方和不变。图~\ref{fig:c3-h2-s}中梯度在两条局部单位基上的投影，正是左侧两个分量。


\begin{figure}[H]\centering
\begin{tikzpicture}[>=Stealth,line cap=round,line join=round]
\fill(0,0)circle(1.5pt)node[below]{$P$};\draw[axis](0,0)--(1.7,1.2)node[right]{$\vct e_s$};\draw[axis](0,0)--(-1.2,1.7)node[left]{$\vct e_t$};\draw[axis](0,0)--(.5,2.4)node[above]{$\nabla f$};\draw[guide](.5,2.4)--(1.46,1.03);
\end{tikzpicture}
\caption{梯度在径向和角向的正交投影，角向尺度因子为$s$。}\label{fig:c3-h2-s}
\end{figure}
\end{solution}

\Needspace{8\baselineskip}
\begin{exercise}[偏导存在但不可微]\label{ex:c3-h3}
函数$f(x,y)=xy/\sqrt{x^2+y^2}$在原点定义为零。求原点的两个偏导，判断可微性，并计算沿$y=x$的单位方向导数。 见图~\ref{fig:c3-h3-q}。
\end{exercise}

\begin{figure}[H]\centering
\begin{tikzpicture}[>=Stealth,line cap=round,line join=round]
\draw[axis](-1.5,0)--(1.7,0)node[right]{$x$};\draw[axis](0,-1.3)--(0,1.7)node[above]{$y$};\draw[curve](-1.2,-1.2)--(1.2,1.2);\node[right]at(1,1){$y=x$};
\end{tikzpicture}
\caption{坐标轴上的函数值为零，斜线上的增量则含一阶项。}\label{fig:c3-h3-q}
\end{figure}
\begin{solution}
原点偏导只需考察两条坐标轴，而可微性要对所有小位移统一成立。图~\ref{fig:c3-h3-q}选取对角线$y=x$，就是为了检验轴向信息是否遗漏了其他方向的一阶变化。方向导数还需说明采用双侧还是单侧极限。


沿坐标轴$f(h,0)=f(0,h)=0$，故$f_x(0,0)=f_y(0,0)=0$。若可微，线性部分应为零，且$f/\sqrt{x^2+y^2}\to0$。然而沿$y=x\ne0$有
\[\frac{f(x,x)}{\sqrt{2x^2}}=\frac12,
\]
所以不可微。图~\ref{fig:c3-h3-s}取单位方向$\vct n=(1,1)/\sqrt2$，则$f(t\vct n)=|t|/2$。按双侧方向导数定义，差商为$\sgn(t)/2$，因此双侧方向导数不存在；若只取$t\to0^+$的单侧增长率，则等于$1/2$。

这里不能套用 $D_{\vct n}f=\nabla f\cdot\vct n$，因为该公式需要可微性。函数本身在原点连续，下面用不等式确认。


先证明连续性：由$2|xy|\le x^2+y^2$，
\[|f(x,y)|\le\tfrac12\sqrt{x^2+y^2}\to0.\]
但可微性要求在减去零线性项后，$f/r\to0$；沿对角线却恒为$1/2$，所以失败发生在线性余项的阶数，而不是函数值本身不连续。

对任意单位方向$\vct n=(a,b)$，有$f(t\vct n)=ab|t|$，双侧差商为$ab\operatorname{sgn}t$。当$ab\ne0$时双侧方向导数不存在；坐标轴方向则为零。图~\ref{fig:c3-h3-s}是$a=b=1/\sqrt2$的尖点截线。存在两个偏导因此既不能保证可微，也不能保证所有双侧方向导数存在。

\begin{figure}[H]\centering
\begin{tikzpicture}[>=Stealth,line cap=round,line join=round]
\begin{scope}[x=1.00000cm,y=1.12500cm]\draw[axis](-2,0)--(2.25,0)node[right]{$t$};\draw[axis](0,-0.2)--(0,1.4)node[above]{$f(t\vct n)$};\draw[curve,domain=-2:2,samples=101]plot(\x,{abs(\x)/2});\node at(1.1,1.3){$|t|/2$};\end{scope}
\end{tikzpicture}
\caption{斜线限制函数有尖点，双侧变化率不同。}\label{fig:c3-h3-s}
\end{figure}
\end{solution}

\Needspace{8\baselineskip}
\begin{exercise}[退化二阶判据]\label{ex:c3-h4}
分类讨论$f=x^4+y^4$、$g=x^4-y^4$、$h=x^2y^2$在原点的极值性质，说明Hessian行列式为零时为何需要继续分析。 见图~\ref{fig:c3-h4-q}。
\end{exercise}

\begin{figure}[H]\centering
\begin{tikzpicture}[>=Stealth,line cap=round,line join=round]
\draw[axis](-1.5,0)--(1.5,0)node[right]{$x$};\draw[axis](0,-1.3)--(0,1.5)node[above]{$y$};\draw[curve](-1,-1)--(1,1);\node[right]at(.8,.8){比较不同路径};
\end{tikzpicture}
\caption{原点的二阶导数都不足以单独决定性质。}\label{fig:c3-h4-q}
\end{figure}
\begin{solution}
三个函数在原点都没有非零二次项，所以Hessian判据无法区分它们。应直接比较足够小邻域内函数值的符号。图~\ref{fig:c3-h4-q}提示沿不同路径观察，但严格极小值还需覆盖所有非零小位移，而非只看有限几条直线。


先实际求导：$\nabla f=(4x^3,4y^3)$、$\nabla g=(4x^3,-4y^3)$、$\nabla h=(2xy^2,2x^2y)$，三个 Hessian 分别为
\[
H_f=\begin{pmatrix}12x^2&0\\0&12y^2\end{pmatrix},\quad
H_g=\begin{pmatrix}12x^2&0\\0&-12y^2\end{pmatrix},\quad
H_h=\begin{pmatrix}2y^2&4xy\\4xy&2x^2\end{pmatrix}.
\]
在原点它们全为零，但函数值符号不同。第一式 $f\ge0$，只有原点为零，是严格全局极小值。第二式沿 $x$ 轴为正、沿 $y$ 轴为负，因此原点是鞍点。第三式 $h\ge0$，原点是非严格全局极小值；两条坐标轴上都为零，所以不是孤立极小点。

图~\ref{fig:c3-h4-s}把第三式的零值集合画成两条轴线。这三个函数的二阶项均为零，极值性质须由四阶项判断。若高阶项也含正负混合，还需寻找能显露其符号的路径。


判断严格极小时还可给出与距离一致的下界：
\[x^4+y^4\ge\tfrac12(x^2+y^2)^2,\]
所以$f$沿任何非零方向离开原点都会增大。对$g$，两条轴上的值分别是$x^4$与$-y^4$，任意小邻域都包含正、负值。对$h$，沿$y=0$始终取零，而沿$y=x\ne0$取正值。解答图因此将零值集合与正值区域区分开来，非严格极小值并不表示函数在邻域内恒定。

若把第三式用作机械势能，沿图~\ref{fig:c3-h4-s}中的零轴没有回复力，切向初速度可引起自由漂移。非严格极小值因而不保证动力学中的完全束缚。

\begin{figure}[H]\centering
\begin{tikzpicture}[>=Stealth,line cap=round,line join=round]
\draw[axis](-1.7,0)--(1.9,0)node[right]{$x$};\draw[axis](0,-1.5)--(0,1.8)node[above]{$y$};\draw[line width=1.3pt](-1.5,0)--(1.5,0) (0,-1.3)--(0,1.5);\node at(2.8,1){$x^2y^2=0$};\node at(2.8,.4){非孤立的最小值集合};
\end{tikzpicture}
\caption{第三式在两条坐标轴上同时取到零值。}\label{fig:c3-h4-s}
\end{figure}
\end{solution}

\Needspace{8\baselineskip}
\begin{exercise}[椭圆约束下的线性极值]\label{ex:c3-h5}
在$x^2/a^2+y^2/b^2=1$上求$ux+vy$的最大最小值，其中$a,b>0$且$(u,v)\ne(0,0)$。 见图~\ref{fig:c3-h5-q}。
\end{exercise}

\begin{figure}[H]\centering
\begin{tikzpicture}[>=Stealth,line cap=round,line join=round]
\draw[curve](0,0)ellipse(2 and 1);\draw[axis](-2.4,0)--(2.6,0)node[right]{$x$};\draw[axis](0,-1.4)--(0,1.9)node[above]{$y$};
\draw[dashed](1.05,{sqrt(5)-1.05})--(2.6,{sqrt(5)-2.6});
\fill({4/sqrt(5)},{1/sqrt(5)})circle(1.5pt)node[above right]{$P_+$};
\node[below]at(2,0){$a$};\node[left]at(0,1){$b$};
\end{tikzpicture}
\caption{图取$a=2,b=1,u=v=1$，虚线$x+y=\sqrt5$恰与椭圆相切，切点为最大值点$P_+$。}\label{fig:c3-h5-q}
\end{figure}
\begin{solution}
椭圆各方向的长度尺度不同。先用$x=aX,y=bY$把可行域化成单位圆，目标函数的系数同时变成$(au,bv)$，再应用内积界。图~\ref{fig:c3-h5-q}中的支撑线在极值处相切，但其法向通常不沿椭圆的径向。


令$X=x/a,Y=y/b$，约束化为$X^2+Y^2=1$，目标为$auX+bvY$。由Cauchy不等式，
\[|ux+vy|\le\sqrt{a^2u^2+b^2v^2}.
\]
等号要求单位向量$(X,Y)=\pm(au,bv)/\sqrt{a^2u^2+b^2v^2}$。变回$x=aX,y=bY$时，分子再各乘一次$a,b$，因此
\[\boxed{f_{\max,\min}=\pm\sqrt{a^2u^2+b^2v^2},\quad
(x,y)=\pm\frac{(a^2u,b^2v)}{\sqrt{a^2u^2+b^2v^2}}}.\]
图~\ref{fig:c3-h5-s}中的法向是$(x/a^2,y/b^2)$，它与$(u,v)$平行，与乘子方程一致。极值点的径向$(x,y)$通常不平行于目标方向，只有圆或特殊轴向情形才重合。闭椭圆保证最大最小值均取得。



图~\ref{fig:c3-h5-s}取$a=2,b=1,u=v=1$，最大点为$(4/\sqrt5,1/\sqrt5)$，其和为$\sqrt5$，且$x^2/4+y^2=1$。法向$(x/a^2,y/b^2)=(1,1)/\sqrt5$与目标方向一致。这个代入也验证了题图的虚线确为相切支撑线，而非仅与曲线相交的一条等值线。

\begin{figure}[H]\centering
\begin{tikzpicture}[>=Stealth,line cap=round,line join=round]
\draw[curve](0,0)ellipse(2 and 1);\coordinate(P)at({4/sqrt(5)},{1/sqrt(5)});
\coordinate(M)at({-4/sqrt(5)},{-1/sqrt(5)});
\fill(P)circle(1.5pt)node[below right]{$P_+$};\fill(M)circle(1.5pt)node[below left]{$P_-$};
\draw[axis](0,0)--(P)node[midway,below]{$\vct r$};
\draw[axis](P)--({4/sqrt(5)+.8},{1/sqrt(5)+.8})node[above]{$(u,v)$};
\draw[dashed]({4/sqrt(5)-.8},{1/sqrt(5)+.8})--({4/sqrt(5)+.6},{1/sqrt(5)-.6});
\end{tikzpicture}
\caption{最大、最小点关于原点对称；图示参数与题图相同。切线法向与$(u,v)$平行，径向通常不与它平行。}\label{fig:c3-h5-s}
\end{figure}
\end{solution}

\Needspace{8\baselineskip}
\begin{exercise}[约束梯度失效的点]\label{ex:c3-h6}
在约束$y^2=x^3$上求$f=x$的最小值。为什么直接使用$\nabla f=\lambda\nabla(y^2-x^3)$会找不到最小点？ 见图~\ref{fig:c3-h6-q}。
\end{exercise}

\begin{figure}[H]\centering
\begin{tikzpicture}[>=Stealth,line cap=round,line join=round]
\draw[axis](-.4,0)--(2.6,0)node[right]{$x$};\draw[axis](0,-2)--(0,2)node[above]{$y$};\draw[curve,domain=0:1.5]plot(\x,{(\x)^1.5});\draw[curve,domain=0:1.5]plot(\x,{-(\x)^1.5});
\end{tikzpicture}
\caption{半三次抛物线的尖点不是正规约束点。}\label{fig:c3-h6-q}
\end{figure}
\begin{solution}
图~\ref{fig:c3-h6-q}中的约束为$y^2=x^3$。由$y^2\ge0$知$x\ge0$，且$x=0$强制$y=0$，所以
$\boxed{f_{\min}=0}$，唯一取得于原点。

设$g=y^2-x^3$，乘子方程及约束为
\[
 1=-3\lambda x^2,\qquad 0=2\lambda y,\qquad y^2=x^3.
\]
第一式要求$\lambda\ne0$，第二式便给$y=0$，约束再给$x=0$，代回第一式矛盾。因此乘子方程没有解。原因是原点处
$\nabla g=(0,0)$，不满足乘子法所需的约束正规条件。
\begin{figure}[H]\centering
\begin{tikzpicture}[>=Stealth,line cap=round,line join=round]
\begin{scope}[x=1.00000cm,y=1.00000cm]\draw[axis](-1.6,0)--(1.85,0)node[right]{$t$};\draw[axis](0,0)--(0,3)node[above]{$f$};\draw[curve,domain=-1.6:1.6,samples=101]plot(\x,{(\x)^2});\fill(0,0)circle(1.5pt)node[below]{$t=0$};\end{scope}
\end{tikzpicture}
\caption{用完整参数化消去约束后，尖点对应一元函数的最小点。}\label{fig:c3-h6-s}
\end{figure}
为完整描述可行曲线，取$x=t^2,y=t^3$，$t\in\mathbb R$。
代入确有$y^2=x^3$；反过来，约束上的非零点有$x>0$，令$t=y/x$，便得
$t^2=y^2/x^2=x$、$t^3=tx=y$，原点则由$t=0$给出。
其中$t>0$对应上支，$t<0$对应下支。

目标函数化为$f(t)=t^2$，有$f'(t)=2t$、$f''(t)=2$。
图~\ref{fig:c3-h6-s}画的是这个参数化后的目标函数，并非原约束曲线。
全局极值讨论须单列$\nabla g=0$的点；乘子方程无解不能否定这些点上存在极值。
\end{solution}

\Needspace{8\baselineskip}


\Needspace{8\baselineskip}
\begin{exercise}[径向场的散度与奇点]\label{ex:c3-h8}
在$r>0$处求三维场$\vct A=C r^n\vct e_r$的散度和旋度。特别讨论$n=-2$时局部散度为零与包围原点的通量非零是否矛盾。 见图~\ref{fig:c3-h8-q}。
\end{exercise}

\begin{figure}[H]\centering
\begin{tikzpicture}[>=Stealth,line cap=round,line join=round]
\draw[curve](0,0)circle(1.1);\foreach\a in{0,45,...,315}{\draw[axis](\a:1.1)--(\a:1.8);}\node at(0,0){$O$};
\end{tikzpicture}
\caption{径向场穿过球面，球面法向与径向一致。}\label{fig:c3-h8-q}
\end{figure}
\begin{solution}
场只有径向分量且只随球心距离变化，使用球坐标能把三维导数化成一维。图~\ref{fig:c3-h8-q}中的球面与场方向正交，通量只需场强乘面积；原点是否属于场域则会影响能否使用实心球的散度定理。


球坐标散度公式在本题化为
\[
\nabla\cdot\vct A=\frac1{r^2}\frac{\dd}{\dd r}(r^2Cr^n)
=\frac C{r^2}(n+2)r^{n+1}=C(n+2)r^{n-1}.
\]
角向分量为零且无角依赖，三个旋度分量均为零，因此
\[\boxed{\nabla\cdot\vct A=C(n+2)r^{n-1},\qquad
\nabla\times\vct A=0\quad(r>0)}.\]
旋度为零也可从势函数确认：对 $n\ne-1$，径向原函数为 $\int Cr^n\dd r=Cr^{n+1}/(n+1)+C_0$；$n=-1$ 时为 $C\ln r+C_0$。它们的梯度均为 $\vct A$。当 $n=-2$，任意半径球面通量为 $4\pi C$，而原点以外散度确实为零。

图~\ref{fig:c3-h8-s}把原点挖去后，球壳内散度积分为零；外表面通量$4\pi C$与内表面通量$-4\pi C$抵消。直接对包含原点的实心球使用光滑场版本的散度定理，违反了正则性条件。若采用分布语言，可将源写为$\nabla\cdot(C\vct e_r/r^2)=4\pi C\delta^{(3)}(\vct r)$；竞赛计算中明确挖去奇点已足够。


球壳半径$\epsilon\le r\le R$上的直接积分为
\[\begin{aligned}
\int_\epsilon^R C(n+2)r^{n-1}\,4\pi r^2\dd r
&=4\pi C(n+2)\left[\frac{r^{n+2}}{n+2}\right]_\epsilon^R\\
&=4\pi C(R^{n+2}-\epsilon^{n+2})\quad(n\ne-2).
\end{aligned}\]
右侧恰为外球面通量减去内球面的径向通量。$n=-2$时应直接代入原散度式，体积分为零，两面通量相等；不能从积分后的式子错误地产生$0/0$。解答图的内边界法向指向球心，这是内面通量带负号的几何来源。

半径 $R$ 球面的通量为场的径向分量乘球面积，即 $CR^n\cdot4\pi R^2=4\pi CR^{n+2}$。


\begin{figure}[H]\centering
\begin{tikzpicture}[>=Stealth,line cap=round,line join=round]
\fill[gray!12](0,0)circle(1.6);\fill[white](0,0)circle(.6);\draw[curve](0,0)circle(1.6);\draw[curve](0,0)circle(.6);\draw[axis](1.6,0)--(2.3,0)node[right]{$\vct n_{\rm out}$};\draw[axis](.6,0)--(.1,0)node[above]{$\vct n_{\rm in}$};\node at(0,-2){内边界法向指向被挖去的区域};
\end{tikzpicture}
\caption{球壳的内边界外法向为$-\vct e_r$，其通量必须计入。}\label{fig:c3-h8-s}
\end{figure}
\end{solution}

\Needspace{8\baselineskip}
\begin{exercise}[刚体转动场]\label{ex:c3-h9}
速度场为$\vct v=\vct\Omega\times\vct r$，其中$\vct\Omega$恒定。取$\vct\Omega=\Omega\vct e_z$求散度、旋度，并求半径$R$圆周上的环量。 见图~\ref{fig:c3-h9-q}。
\end{exercise}

\begin{figure}[H]\centering
\begin{tikzpicture}[>=Stealth,line cap=round,line join=round]
\draw[curve](0,0)ellipse(1.6 and .7);\draw[axis](0,0)--(0,2)node[above]{$\vct\Omega$};\draw[axis](1.6,0)--(1.6,.8)node[right]{$\vct v$};
\end{tikzpicture}
\caption{圆周切向速度随轴距成正比。}\label{fig:c3-h9-q}
\end{figure}
\begin{solution}
绕$z$轴匀转时，速度与轴距成正比并沿方位角方向。图~\ref{fig:c3-h9-q}的圆面可视作刚体的一块截面。散度检验体积是否局部改变，旋度检验环量密度，二者在此给出不同的信息。


直角坐标分量为 $(-\Omega y,\Omega x,0)$。散度的三个导数分别为 $\partial_x(-\Omega y)=0$、$\partial_y(\Omega x)=0$、$\partial_z0=0$；旋度 $z$ 分量为 $\partial_x(\Omega x)-\partial_y(-\Omega y)=\Omega-(-\Omega)=2\Omega$，其余分量为零。因此
\[\boxed{\nabla\cdot\vct v=0,\qquad\nabla\times\vct v=2\Omega\vct e_z=2\vct\Omega}.
\]
图~\ref{fig:c3-h9-s}中从正 $z$ 方向看取逆时针回路，线元为 $R\vct e_\varphi\dd\varphi$，速度为 $\Omega R\vct e_\varphi$，带符号的点积是 $\Omega R^2\dd\varphi$。所以
\[
\oint\vct v\cdot\dd\vct l=\int_0^{2\pi}\Omega R^2\dd\varphi
=\Omega R^2[\varphi]_0^{2\pi},
\]
得到
\[\boxed{\oint\vct v\cdot\dd\vct l=2\pi\Omega R^2}.
\]
除以圆面积得到$2\Omega$，与旋度法向分量一致。旋度是局部环量密度，它对刚体转动给出角速度的两倍；不能直接把$\nabla\times\vct v$称为刚体角速度。场不可压缩并不意味着没有转动。


面积分法给出 $2\Omega\cdot\pi R^2$，与直接环量积分相同。因子二来自小回路两组边的速度差，并非坐标选择造成的额外倍数。

\begin{figure}[H]\centering
\begin{tikzpicture}[>=Stealth,line cap=round,line join=round]
\draw[curve](0,0)circle(1.4);\draw[axis](1.4,-.2)--(1.4,.3)node[right]{$\dd\vct l$};\draw[axis](0,0)--(1.4,0)node[midway,below]{$R$};\draw[axis](0,1.4)--(-.7,1.4)node[above]{$\vct v$};\node at(0,0){$+z\ \odot$};
\end{tikzpicture}
\caption{正向环路与正$z$法向按右手规则配对。}\label{fig:c3-h9-s}
\end{figure}
\end{solution}

\Needspace{8\baselineskip}
\begin{exercise}[有孔区域中的无旋场]\label{ex:c3-h10}
在去掉$z$轴的空间中，$\vct A=K\vct e_\varphi/r$，$r$为柱坐标轴距。求旋度和绕轴一周的环量，讨论单值势函数是否存在。 见图~\ref{fig:c3-h10-q}。
\end{exercise}

\begin{figure}[H]\centering
\begin{tikzpicture}[>=Stealth,line cap=round,line join=round]
\draw[curve](0,0)circle(1.5);\fill(0,0)circle(2pt);\draw[axis](1.5,-.2)--(1.5,.3);\node at(0,-2){轴线从定义域中挖去};
\end{tikzpicture}
\caption{回路不能在场的定义域内连续收缩成一点。}\label{fig:c3-h10-q}
\end{figure}
\begin{solution}
场在轴线以外处处光滑，局部旋度可直接由柱坐标求出。但定义域已排除整条轴，绕轴闭路不能在域内收缩成一点。图~\ref{fig:c3-h10-q}明确这个孔洞，正是局部无旋不足以推出全域单值势的原因。


柱坐标旋度的$z$分量为$r^{-1}\partial_r(rA_\varphi)=r^{-1}\partial_rK=0$，其他分量也为零。但圆周环量为
\[\boxed{\oint\vct A\cdot\dd\vct l=\int_0^{2\pi}(K/r)r\dd\varphi
=[K\varphi]_0^{2\pi}=2\pi K}.
\]
若存在全域单值势，闭合回路积分必为零，故$K\ne0$时不存在。局部可写$\vct A=\nabla(K\varphi)$，然而绕轴一周后$\varphi$增加$2\pi$，局部势不能在整个有孔区域拼成单值函数。



绕轴半径为$R$时，场强$K/R$与线元长度$R\dd\varphi$相乘，轴距恰好抵消，得到与回路半径无关的$2\pi K$。若$K=0$，场恒零，当然存在常势；非零结论应明确以$K\ne0$为条件。

图~\ref{fig:c3-h10-s}沿一条半平面作切口后，可以选$0<\varphi<2\pi$的单值角函数，其梯度为$\vct e_\varphi/r$。切口两侧角度极限相差$2\pi$，势值相差$2\pi K$，所以不能在未切开的全域连续拼合。以跨过轴线的圆盘直接应用Stokes公式，会让曲面穿过场的奇点，违背其前提。

\begin{figure}[H]\centering
\begin{tikzpicture}[>=Stealth,line cap=round,line join=round]
\draw[curve](0,0)circle(1.5);\draw[thick](0,0)--(2,0);\node[above]at(1,.05){$\varphi=0^+$};\node[below]at(1,-.05){$\varphi=2\pi^-$};\draw[axis](-1.5,.2)--(-1.5,-.3);\node at(0,-2){局部势$K\varphi$在切口两侧相差$2\pi K$};
\end{tikzpicture}
\caption{切开区域可定义局部角度，但势的跳变记录了绕轴环量。}\label{fig:c3-h10-s}
\end{figure}
\end{solution}

\Needspace{8\baselineskip}
\begin{exercise}[径向拉普拉斯方程]\label{ex:c3-h11}
分别求柱对称$f=f(r)$与球对称$f=f(r)$在无源区域满足$\nabla^2f=0$的全部形式，并说明两种$r$的区别。 见图~\ref{fig:c3-h11-q}。
\end{exercise}

\begin{figure}[H]\centering
\begin{tikzpicture}[>=Stealth,line cap=round,line join=round]
\draw[curve](-1,-1)--(-1,1);\draw[curve](1,-1)--(1,1);\draw[curve](0,1)ellipse(1 and .3);\draw[curve](0,-1)ellipse(1 and .3);\begin{scope}[shift={(4,0)}]\draw[curve](0,0)circle(1);\draw[dashed](0,0)ellipse(1 and .3);\end{scope}
\end{tikzpicture}
\caption{柱对称与球对称的通量面积分别随$r$与$r^2$增长。}\label{fig:c3-h11-q}
\end{figure}
\begin{solution}
“柱对称”和“球对称”的径向距离不同，因此无源方程中的几何面积因子也不同。图~\ref{fig:c3-h11-q}分别给出柱面与球面。把各自的对称性代入Laplace算子后，径向积分将分别产生$\ln r$与$1/r$。


柱对称且与 $z,\varphi$ 无关时，$r$ 是到轴的距离，方程为 $(rf')'=0$。第一次积分给 $rf'=C_1$，故 $f'=C_1/r$；第二次积分为 $f=C_1\int r^{-1}\dd r=C_1\ln r+C_2$，即
\[\boxed{f=A\ln r+B\qquad(r>0)}.
\]
球对称时 $r$ 为到球心的距离，方程为 $(r^2f')'=0$。第一次积分给 $r^2f'=C_1$，第二次积分为 $f=C_1\int r^{-2}\dd r=-C_1/r+C_2$。将 $-C_1$ 改记为 $A$，得到
\[\boxed{f=A/r+B\qquad(r>0)}.
\]
图~\ref{fig:c3-h11-s}中等势面与等通量面相配：柱面单位长度面积$2\pi r$，球面面积$4\pi r^2$，因此梯度分别按$1/r$、$1/r^2$变化。常数的正负可吸收入$A$。

若要求包含原点或轴线的区域内光滑有限，奇异项系数必须为零；若存在点源或线源，奇异项正是物理模型所需。边界条件决定$A,B$，仅有无源微分方程不能唯一确定势。


设同一无源柱壳或球壳的两边界分别为 $r_1,r_2$，势值为 $f_1,f_2$，其中 $0<r_1<r_2$。柱情形两式相减，$f_2-f_1=A\ln(r_2/r_1)$，所以 $A=(f_2-f_1)/\ln(r_2/r_1)$、$B=f_1-A\ln r_1$；球情形则有 $f_2-f_1=A(r_2^{-1}-r_1^{-1})$，所以 $A=(f_2-f_1)/(r_2^{-1}-r_1^{-1})$、$B=f_1-A/r_1$。代入得到
\[f_{\rm cyl}(r)=f_1+(f_2-f_1)\frac{\ln(r/r_1)}{\ln(r_2/r_1)},\]
\[f_{\rm sph}(r)=f_1+(f_2-f_1)\frac{r^{-1}-r_1^{-1}}{r_2^{-1}-r_1^{-1}}.\]
图中柱面与球面的面积增长率不同，决定了两种不同的径向插值。这里仅在不含源的壳层内求解，内边界的源可以存在于被排除的区域。


图~\ref{fig:c3-h11-s}的曲线展示的是势函数，而通量守恒控制的是势梯度乘相应面积。若柱区域一直延伸到无穷远，非零$\ln r$不能像$1/r$那样取有限零势，所以实际柱对称势常以有限参考半径定义势差。相同的无源局部方程在不同边界区域上会有不同的可用常数选择。

\begin{figure}[H]\centering
\begin{tikzpicture}[>=Stealth,line cap=round,line join=round]
\begin{scope}[x=1.00000cm,y=1.00000cm]\draw[axis](0,0)--(3.25,0)node[right]{$r$};\draw[axis](0,-1.5)--(0,2)node[above]{$f$};\draw[curve,domain=0.25:3,samples=101]plot(\x,{ln(\x)});\draw[dashed,domain=.5:3,samples=101]plot(\x,{1/\x});\node at(2.3,1.8){实线$\ln r$，虚线$1/r$};\end{scope}
\end{tikzpicture}
\caption{两种径向势由不同的几何扩散面积产生。}\label{fig:c3-h11-s}
\end{figure}
\end{solution}

\Needspace{8\baselineskip}
\begin{exercise}[偏心电场中的偶极子受力]\label{ex:c3-h12}
小永久电偶极矩$\vct p$固定不转，置于缓变静电场$\vct E(\vct r)$中。由两点电荷受力展开求合力，并在场无旋处证明$\vct F=\nabla(\vct p\cdot\vct E)$。 见图~\ref{fig:c3-h12-q}。
\end{exercise}

\begin{figure}[H]\centering
\begin{tikzpicture}[>=Stealth,line cap=round,line join=round]
\fill(0,0)circle(2pt)node[below right]{$-q$};\fill(1.6,.8)circle(2pt)node[above left]{$+q$};\draw[axis](0,0)--(1.6,.8)node[midway,above]{$\vct\ell$};\draw[axis](0,0)--(-.5,-.6)node[below left]{$-q\vct E_-$};\draw[axis](1.6,.8)--(2.3,1.7)node[above right]{$q\vct E_+$};
\end{tikzpicture}
\caption{箭头表示两电荷的受力；负电荷受力与当地外场反向，两力因场的非均匀性而不能完全抵消。}\label{fig:c3-h12-q}
\end{figure}
\begin{solution}
图~\ref{fig:c3-h12-q}中$\vct\ell$由负电荷指向正电荷，以偶极中心$\vct r$为展开点，两电荷分别位于$\vct r\pm\vct\ell/2$。假定外场在偶极尺度内足够光滑，逐分量展开：
\[E_i(\vct r\pm\vct\ell/2)=E_i(\vct r)
\pm\frac12\ell_j\partial_jE_i
+\frac18\ell_j\ell_k\partial_j\partial_kE_i+O(\ell^3\|\nabla^3\vct E\|).
\]
重复指标表示求和。两场值相减时常数项与二次项抵消；乘上$q$，令$\vct p=q\vct\ell$，得到
\[F_i=p_j\partial_jE_i+O(q\ell^3\|\nabla^3\vct E\|),\qquad
\vct F=(\vct p\cdot\nabla)\vct E
\]
（后一式取理想偶极近似）。图~\ref{fig:c3-h12-s}中的差分沿偶极轴作用于每个场分量，其方向由$\vct\ell$固定。

本题的$\vct p$在平移求导时保持不变，故
\[[\nabla(\vct p\cdot\vct E)]_i=p_j\partial_iE_j.
\]
静电场无旋给出$\partial_iE_j=\partial_jE_i$，从而
\[\boxed{\vct F=(\vct p\cdot\nabla)\vct E
=\nabla(\vct p\cdot\vct E)}.\]
由永久偶极的势能$U=-\vct p\cdot\vct E$取负梯度，也得到同一平移力。若解除转动约束，角度变化还给出力矩$\vct\tau=\vct p\times\vct E$。对于随场改变的感生偶极，应先计入极化所需能量，不能把位置相关的$\vct p$直接代入永久偶极的能量后照常求导。

\begin{figure}[H]\centering
\begin{tikzpicture}[>=Stealth,line cap=round,line join=round]
\node at(0,0){$\vct E_-$};\node at(3,0){$\vct E_+$};\draw[axis](.7,0)--(2.3,0)node[midway,above]{$\vct\ell\cdot\nabla$};\draw[axis](4,0)--(6,.5)node[midway,above]{$q(\vct E_+-\vct E_-)$};\node at(3,-1){$F_i=p_j\partial_jE_i$};
\end{tikzpicture}
\caption{逐分量展开后再利用无旋条件，避免把标量梯度与向量方向导数混淆。}\label{fig:c3-h12-s}
\end{figure}
\end{solution}

\Needspace{8\baselineskip}
\begin{exercise}[定温与绝热偏导的区别]\label{ex:c3-h13}
物质的量$n$固定的理想气体满足$pV=nRT$。考虑准静态可逆绝热过程，摩尔热容近似为常数，$\kappa=C_p/C_V>1$，过程关系为$pV^\kappa=C$。求$(\partial p/\partial V)_T$和沿绝热过程的$\dd p/\dd V$，并说明为什么不能在第一个偏导中同时令$T$按绝热关系变化。 见图~\ref{fig:c3-h13-q}。
\end{exercise}

\begin{figure}[H]\centering
\begin{tikzpicture}[>=Stealth,line cap=round,line join=round]
\begin{scope}[x=1.30233cm,y=1.00000cm]\draw[axis](0,0)--(2.75,0)node[right]{$V/V_0$};\draw[axis](0,0)--(0,2)node[above]{$p/p_0$};\draw[curve,domain=0.6:2.5,samples=101]plot(\x,{1/\x});\draw[dashed,domain=.6:2.5]plot(\x,{(\x)^(-1.4)});\end{scope}
\end{tikzpicture}
\caption{实线为定温过程，虚线为$\kappa=1.4$的可逆绝热过程；两条曲线经过同一状态。}\label{fig:c3-h13-q}
\end{figure}
\begin{solution}
图~\ref{fig:c3-h13-q}的两条过程曲线经过同一状态，但求导时保持的量不同。定温时$p=nRT/V$中的$nRT$为常数，所以
\[\boxed{\left(\frac{\partial p}{\partial V}\right)_T=-\frac{nRT}{V^2}=-\frac pV}.\]

对题设的准静态可逆绝热过程，第一定律为$nC_V\dd T=-p\dd V$。结合$pV=nRT$以及$C_p-C_V=R$，有
\[\frac{\dd T}{T}=-(\kappa-1)\frac{\dd V}{V}.
\]
从共同状态$(V_0,T_0)$积分，常热容近似使$\kappa$可视为常数，因而
\[
\int_{T_0}^{T}\frac{\dd u}{u}
=-(\kappa-1)\int_{V_0}^{V}\frac{\dd v}{v},\qquad
\ln\frac{T}{T_0}=-(\kappa-1)\ln\frac{V}{V_0}.
\]
指数化得$TV^{\kappa-1}=T_0V_0^{\kappa-1}$，与状态方程联立即得$pV^\kappa=C$。沿此曲线对$V$求导，乘积法则给出
\[V^\kappa\frac{\dd p}{\dd V}+\kappa pV^{\kappa-1}=0,\]
除以$V^\kappa$得到
\[\boxed{\frac{\dd p}{\dd V}=-\kappa\frac pV}.\]

由全微分也能看出区别：
\[\dd p=-\frac pV\dd V+\frac pT\dd T.\]
定温令$\dd T=0$；可逆绝热则将上述$\dd T/\dd V$代入，两项合成$-\kappa p/V$。图~\ref{fig:c3-h13-s}比较同一点$P$的切线，绝热切线斜率绝对值是定温切线的$\kappa$倍。在定温偏导中再让温度随绝热关系变化，实际已换成另一条路径的全导数。

绝热只要求没有热量交换，本身不保证$pV^\kappa$不变。例如理想气体向真空自由膨胀也可绝热，但过程不可逆，初、末状态不遵从本题的可逆绝热幂律。若热容显著随温度变化，应积分$C_V(T)\dd T/T=-R\dd V/V$，不能预先使用常指数形式。

\begin{figure}[H]\centering
\begin{tikzpicture}[>=Stealth,line cap=round,line join=round]
\draw[axis](0,0)--(3,0)node[right]{$V$};\draw[axis](0,0)--(0,2.7)node[above]{$p$};\fill(1.5,1.5)circle(1.5pt)node[above right]{$P$};\draw[curve](.6,2.1)--(2.4,.9)node[above right]{定温};\draw[dashed](.9,2.06)--(2.1,.94)node[below right]{可逆绝热};
\end{tikzpicture}
\caption{在同一点比较切线，可逆绝热斜率的绝对值为定温斜率的$\kappa$倍；图取$\kappa=1.4$。}\label{fig:c3-h13-s}
\end{figure}
\end{solution}

\Needspace{8\baselineskip}
\begin{exercise}[球面上线性势的二阶判别]\label{ex:c3-h14}
质量$m$的质点被约束在球面$x^2+y^2+z^2=R^2$上，重力势能$U=mgz$。求平衡位置并用约束二阶变分判断稳定性；为什么$H_U=0$并不意味着中性平衡？ 见图~\ref{fig:c3-h14-q}。
\end{exercise}

\begin{figure}[H]\centering
\begin{tikzpicture}[>=Stealth,line cap=round,line join=round]
\draw[curve](0,0)circle(1.5);\draw[axis](0,-1.9)--(0,2)node[above]{$z$};\fill(0,1.5)circle(2pt)node[right]{$N$};\fill(0,-1.5)circle(2pt)node[right]{$S$};\draw[axis](2,.8)--(2,-.3)node[midway,right]{$mg$};
\end{tikzpicture}
\caption{球面的曲率把线性高度势转化为切向的二阶回复或偏离。}\label{fig:c3-h14-q}
\end{figure}
\begin{solution}
质点沿球面移动时，切向位移会伴随二阶高度变化。图~\ref{fig:c3-h14-q}标出南北两极；用乘子方程找出平衡后，在各点的切平面上计算势能二次变化。


记约束函数$h=x^2+y^2+z^2-R^2$。其梯度在球面上非零，乘子方程为
\[0=2\lambda x,\qquad 0=2\lambda y,\qquad mg=2\lambda z.\]
最后一式使$\lambda\ne0$，前两式便给出$x=y=0$，再由约束得$z=\pm R$。北极$N$处$\lambda=mg/(2R)$，南极$S$处$\lambda=-mg/(2R)$，没有其他平衡位置。

在两极，切向位移满足$\dd z=0$。由于$H_U=0$、$H_h=2I$，拉格朗日Hessian及其切向二次型为
\[H_U-\lambda H_h=-2\lambda I,\qquad
\delta^2U\big|_{\rm tangent}=-2\lambda\bigl[(\dd x)^2+(\dd y)^2\bigr].\]
故北极切向二次型负定，南极正定。

图~\ref{fig:c3-h14-s}在南极附近直接展开$z=-\sqrt{R^2-x^2-y^2}=-R+(x^2+y^2)/(2R)+O(r_\perp^4/R^3)$，于是
\[U=-mgR+\frac{mg}{2R}(x^2+y^2)+\cdots.
\]
南极为势能严格极小点，北极的二阶项取负号，故不稳定。$H_U=0$只描述不受约束的直线位移，球面上的实际可行位移还带来二阶高度变化，必须计入约束曲率。


在南极附近，速度的竖直分量为二阶小量，二次动能为$T_2=m(\dot x^2+\dot y^2)/2$。与图中展开的二次势能相配，线性方程为
\[\ddot x+(g/R)x=0,\qquad \ddot y+(g/R)y=0.\]
代入$x\propto e^{st}$得$s^2+g/R=0$，两根为$s=\pm\ii\sqrt{g/R}$，故两个切向模态频率均为$\sqrt{g/R}$。北极附近$z=R-(x^2+y^2)/(2R)+\cdots$，线性方程改为$\ddot x-(g/R)x=0$及相同的$y$式，其特征根为$s=\pm\sqrt{g/R}$，含增长指数。图~\ref{fig:c3-h14-s}所示的正高度增量正是南极产生切向回复力的原因。

\begin{figure}[H]\centering
\begin{tikzpicture}[>=Stealth,line cap=round,line join=round]
\draw[curve,domain=-1.3:1.3]plot(\x,{-sqrt(2.25-(\x)^2)});\draw[axis](-1.6,-1.5)--(1.7,-1.5)node[right]{$x$};\draw[guide](.8,-1.5)--(.8,-1.268858);\node[right]at(.8,-1.15){$\Delta z\simeq x^2/(2R)$};\fill(0,-1.5)circle(2pt)node[below]{$S$};
\end{tikzpicture}
\caption{切向的一阶位移产生正的二阶高度变化，形成局部势阱。}\label{fig:c3-h14-s}
\end{figure}
\end{solution}

\Needspace{8\baselineskip}
\begin{exercise}[距离函数的梯度与Hessian]\label{ex:c3-h15}
对$\Phi(\vct r)=1/|\vct r|$，$r>0$，求梯度与二阶偏导矩阵，并证明其迹为零。 见图~\ref{fig:c3-h15-q}。
\end{exercise}

\begin{figure}[H]\centering
\begin{tikzpicture}[>=Stealth,line cap=round,line join=round]
\draw[axis](0,0)--(2.8,0)node[right]{$x$};\draw[axis](0,0)--(0,2.4)node[above]{$y$};\draw[axis](0,0)--(1.8,1.2)node[midway,above]{$\vct r$};\draw[curve](0,0)circle(1.3);
\end{tikzpicture}
\caption{函数只依赖径向距离，但二阶导数还区分径向与横向。}\label{fig:c3-h15-q}
\end{figure}
\begin{solution}
图~\ref{fig:c3-h15-q}中的函数只依赖半径。求Hessian时仍要保留方向的变化：先对$r=\sqrt{x_ix_i}$求导，再对梯度的每一分量作乘积求导。


由$r=(x_kx_k)^{1/2}$，链式法则给出
\[\partial_i r=\frac12(x_kx_k)^{-1/2}(2x_i)=\frac{x_i}{r},\qquad
\partial_i\Phi=-r^{-2}\partial_i r=-x_ir^{-3}.\]
再用乘积法则对$x_j$求导，
\[\partial_j(-x_ir^{-3})=-\delta_{ij}r^{-3}
-x_i(-3r^{-4})\frac{x_j}{r}
=-\frac{\delta_{ij}}{r^3}+\frac{3x_ix_j}{r^5},\]
两项分别来自坐标分量和半径幂次。因此
\[\boxed{\partial_j\partial_i\Phi=\frac{3x_ix_j-r^2\delta_{ij}}{r^5}}.
\]
取迹时$\sum_i3x_i^2=3r^2$与$\sum_i r^2=3r^2$抵消，故$\nabla^2\Phi=0$。图~\ref{fig:c3-h15-s}中沿径向的Hessian特征值为$2/r^3$，任一垂直径向的方向为$-1/r^3$，三个特征值之和为零。

零迹并不等于Hessian为零；场在径向与横向有不同曲率。上述结论只在$r>0$成立，原点的源项仍由包围通量决定，不能凭零迹把点源删除。该矩阵也给出了中心引力潮汐场的拉伸与压缩方向。


将矩阵记为$H=(3\vct e_r\vct e_r^{\mathsf T}-I)/r^3$，可直接验证主方向：
\[H\vct e_r=2\vct e_r/r^3,\qquad
H\vct t=-\vct t/r^3\quad(\vct t\cdot\vct e_r=0).\]
解答图中的两个横向基矢张成垂直于$\vct e_r$的平面，两个横向方向的曲率相同，任取正交横向基矢均可。若用于引力势$-GM/r$，须先乘系数$-GM$，再由$\delta\vct F=-H_U\delta\vct r$求力差；不能把电势函数的曲率符号直接当作加速度差的符号。

\begin{figure}[H]\centering
\begin{tikzpicture}[>=Stealth,line cap=round,line join=round]
\fill(0,0)circle(1.5pt)node[below]{$P$};\draw[axis](0,0)--(2,0)node[right]{$\vct e_r:2/r^3$};\draw[axis](0,0)--(0,1.5)node[above]{$\vct e_1:-1/r^3$};\draw[axis](0,0)--(-1,-1)node[left]{$\vct e_2:-1/r^3$};
\end{tikzpicture}
\caption{一个径向与两个横向的二阶曲率相加为零。}\label{fig:c3-h15-s}
\end{figure}
\end{solution}

\Needspace{8\baselineskip}
\begin{exercise}[球对称Poisson方程]\label{ex:c3-h16}
半径$R$的均匀带电球密度$\rho$，真空介电常数$\varepsilon_0$，电势在无穷远取零。由Poisson方程求球内外电势，并说明球心、表面和无穷远的条件怎样确定常数。 见图~\ref{fig:c3-h16-q}。
\end{exercise}

\begin{figure}[H]\centering
\begin{tikzpicture}[>=Stealth,line cap=round,line join=round]
\fill[gray!10](0,0)circle(1.6);\draw[curve](0,0)circle(1.6);\draw[dashed](0,0)circle(.9);
\draw[axis](0,0)--(2.5,0)node[right]{$r$};
\draw[<->](0,-.25)--(.9,-.25)node[midway,below]{$r<R$};
\node[above]at(1.6,0){$R$};\node at(-.75,.75){$\rho$};\node[below left]at(0,0){$O$};
\end{tikzpicture}
\caption{灰色区域是均匀带电实心球，虚线是球内半径$r$的辅助高斯球；它不是空腔边界。}\label{fig:c3-h16-q}
\end{figure}
\begin{solution}
球对称性把Poisson方程降成径向方程，但球内和球外的源项不同，必须分别积分并在表面匹配。图~\ref{fig:c3-h16-q}的内虚线是辅助高斯球而非空腔，球心没有额外点源，因此球内不能保留$1/r$奇异项。


球内$(r^2\Phi')'=-\rho r^2/\varepsilon_0$。先后两次积分为
\[
r^2\Phi'=-\frac{\rho}{\varepsilon_0}\int r^2\dd r+C
=-\frac{\rho r^3}{3\varepsilon_0}+C,
\]
\[
\Phi_{\rm in}=\int\left(-\frac{\rho r}{3\varepsilon_0}+Cr^{-2}\right)\dd r
=-\frac{\rho r^2}{6\varepsilon_0}-\frac Cr+B.
\]
若$C\ne0$，球心便有$1/r$奇异项，对应额外点电荷。题设只有均匀体密度，球心电势有限，故$C=0$。

球外$(r^2\Phi')'=0$，一次积分$r^2\Phi'=C_{\rm out}$，二次积分为
\[\Phi_{\rm out}=C_{\rm out}\int r^{-2}\dd r+B_{\rm out}
=-\frac{C_{\rm out}}r+B_{\rm out}.\]
令$A=-C_{\rm out}$，无穷远零势给出$B_{\rm out}=0$。

表面没有独立面电荷，法向电场连续，故$\Phi'$连续；电势本身也连续。两条件依次写成
\[
-\frac{\rho R}{3\varepsilon_0}=-\frac{A}{R^2}
\quad\Longrightarrow\quad A=\frac{\rho R^3}{3\varepsilon_0},
\]
\[
-\frac{\rho R^2}{6\varepsilon_0}+B=\frac AR
\quad\Longrightarrow\quad B=\frac{\rho R^2}{3\varepsilon_0}
+\frac{\rho R^2}{6\varepsilon_0}=\frac{\rho R^2}{2\varepsilon_0}.
\]
因此
\[\boxed{\Phi(r)=\begin{cases}\rho(3R^2-r^2)/(6\varepsilon_0),&r\le R,\\
\rho R^3/(3\varepsilon_0r),&r\ge R.\end{cases}}\]
图~\ref{fig:c3-h16-s}中两段电势与斜率在$r=R$相接，球心斜率为零。二阶导数可以跳变，因为体密度在表面发生跳变；不应额外要求所有阶导数连续。


由$E_r=-\Phi'$还能得到
\[E_r=\begin{cases}\rho r/(3\varepsilon_0),&r<R,\\
\rho R^3/(3\varepsilon_0r^2),&r>R.\end{cases}\]
两段在$R$处取相同值。对于题图中球心到半径$r<R$的高斯球，包围电荷$Q_r=4\pi\rho r^3/3$，因而$E_r4\pi r^2=Q_r/\varepsilon_0$，与Poisson积分相同。这里要求连续的是$\Phi$和法向电场；若另加面电荷$\sigma$，第二个匹配条件应改为$E_{r,\rm out}-E_{r,\rm in}=\sigma/\varepsilon_0$。

\begin{figure}[H]\centering
\begin{tikzpicture}[>=Stealth,line cap=round,line join=round]
\begin{scope}[x=2.24000cm,y=1.00000cm]\draw[axis](0,0)--(1.25,0);\draw[axis](0,0)--(0,1.8)node[above]{$\Phi/\Phi(R)$};\draw[curve,domain=0:1,samples=101]plot(\x,{(3-(\x)^2)/2});\draw[curve,domain=1:3]plot(\x,{1/\x});\draw[axis](1,0)--(3.3,0)node[right]{$r/R$};\draw[guide](1,0)--(1,1.5);\node[below]at(1,0){$1$};\end{scope}
\end{tikzpicture}
\caption{球内为抛物线，球外为反比例曲线；连接时匹配电势和法向电场。}\label{fig:c3-h16-s}
\end{figure}
\end{solution}

\Needspace{8\baselineskip}
\begin{exercise}[二维无源区域不能有严格势阱]\label{ex:c3-h17}
平面区域内势$U(x,y)$二阶连续且满足$U_{xx}+U_{yy}=0$。若某驻点Hessian非退化，判断它能否为稳定势能极小点；以$U=a(x^2-y^2)$说明。再讨论Hessian退化时，光滑无源区域内部是否仍能出现严格局部极小值。 见图~\ref{fig:c3-h17-q}。
\end{exercise}

\begin{figure}[H]\centering
\begin{tikzpicture}[>=Stealth,line cap=round,line join=round]
\draw[axis](-1.7,0)--(1.8,0)node[right]{$x$};\draw[axis](0,-1.7)--(0,1.8)node[above]{$y$};\draw[axis](.5,0)--(.1,0);\draw[axis](-.5,0)--(-.1,0);\draw[axis](0,.5)--(0,1.1);\draw[axis](0,-.5)--(0,-1.1);
\end{tikzpicture}
\caption{示意$a>0$时，沿$x$恢复而沿$y$偏离。}\label{fig:c3-h17-q}
\end{figure}
\begin{solution}
在题设非退化驻点，Hessian为实对称矩阵。Laplace条件使其两个实特征值之和为零；非退化使它们都不为零，故一正一负，驻点为鞍点。等价地，写$H=\begin{pmatrix}A&B\\B&-A\end{pmatrix}$，便有$\det H=-A^2-B^2<0$。

图~\ref{fig:c3-h17-q}取$a>0$，$U=a(x^2-y^2)$的力为$(-2ax,2ay)$。沿$x$有恢复力，沿$y$却有发散力；质量$m$的质点在两个方向分别具有频率和增长率$\sqrt{2a/m}$。图~\ref{fig:c3-h17-s}的正负等势线显示原点附近总有升高和降低方向。

Hessian退化时不能继续依靠二阶判据，但可直接利用Laplace方程。取完全包含在光滑无源区域内的圆盘$D_r$，中心为$P$，定义圆周平均
\[M(r)=\frac1{2\pi}\int_0^{2\pi}U(P+r\vct e_r(\theta))\dd\theta,
\quad \vct e_r(\theta)=(\cos\theta,\sin\theta).\]
图~\ref{fig:c3-h17-mean}中的外法向即$\vct e_r(\theta)$，弧长元为$r\dd\theta$。由平面散度定理，
\[M'(r)=\frac1{2\pi}\int_0^{2\pi}\partial_rU\dd\theta
=\frac1{2\pi r}\oint_{\partial D_r}\nabla U\cdot\vct n\dd s
=\frac1{2\pi r}\iint_{D_r}\Delta U\dd A=0.\]
由连续性$M(0)=U(P)$，所以每个足够小的圆周平均都等于中心值。若$P$是严格局部极小点，小圆周上处处有$U>U(P)$，其平均必大于$U(P)$，矛盾。即使只要求局部极小，圆周上$U\ge U(P)$与平均相等也迫使每个小圆周处处取等，因而一个邻域内$U$为常数；这种常值区域没有回复力。

只要势在区域内部光滑且调和，圆盘论证便适用，Hessian退化也不例外。圆盘内存在源或奇异点，以及运动受边界限制时，须按相应条件另行分析。

\begin{figure}[H]\centering
\begin{tikzpicture}[>=Stealth,line cap=round,line join=round]
\draw[axis](-1.75,0)--(1.85,0)node[right]{$x$};\draw[axis](0,-1.75)--(0,1.9)node[above]{$y$};
\foreach\k in{.35,.8}{
\draw[curve,domain=-1.3:1.3,samples=81]plot({sqrt((\x)^2+\k)},\x);
\draw[curve,domain=-1.3:1.3,samples=81]plot({-sqrt((\x)^2+\k)},\x);
\draw[dashed,domain=-1.3:1.3,samples=81]plot(\x,{sqrt((\x)^2+\k)});
\draw[dashed,domain=-1.3:1.3,samples=81]plot(\x,{-sqrt((\x)^2+\k)});}
\node[right,align=left]at(2.5,.7){实线：$U>0$\\[5pt]虚线：$U<0$\\[7pt]$\lambda_1+\lambda_2=0$};
\end{tikzpicture}
\caption{$U=a(x^2-y^2),a>0$的完整正、负等势线。每一条双曲线的两支均保留，原点附近有升高和降低方向。}\label{fig:c3-h17-s}
\end{figure}

\begin{figure}[H]\centering
\begin{tikzpicture}[>=Stealth]
\draw[curve](0,0)circle(1.45);\fill(0,0)circle(1.5pt)node[below left]{$P$};
\draw[guide](0,0)--(1.025,1.025)node[midway,above left]{$r$};
\fill(1.025,1.025)circle(1.5pt);\draw[axis](1.025,1.025)--(1.65,1.65)node[above]{$\vct n$};
\draw[axis](1.45,0)arc(0:28:1.45);\node[right]at(1.48,.42){$\dd s=r\dd\theta$};
\node at(-.55,.5){$D_r$};\node[right,align=left]at(3.4,.3){$\Delta U=0$\\[6pt]$M(r)=U(P)$};
\end{tikzpicture}
\caption{圆盘上的无源条件使外法向梯度总通量为零，圆周平均不随半径变化。}\label{fig:c3-h17-mean}
\end{figure}

\end{solution}

\Needspace{8\baselineskip}
\begin{exercise}[流场中的材料导数]\label{ex:c3-h18}
温度场$T(x,y,t)=T_0+ax+by+ct$，流速场$\vct v=(u,-\lambda y)$，$u,\lambda>0$。一流体微团初位于$(x_0,y_0)$，求它测得的温度随时间变化，并比较固定温度计的读数。 见图~\ref{fig:c3-h18-q}。
\end{exercise}

\begin{figure}[H]\centering
\begin{tikzpicture}[>=Stealth,line cap=round,line join=round]
\draw[axis](0,0)--(4,0)node[right]{$x$};\draw[axis](0,0)--(0,3)node[above]{$y$};\draw[curve,domain=0:3]plot(\x,{2*exp(-.6*\x)});\draw[axis](.6,1.4)--(1.2,1.05);\node[right]at(2,2.3){流体轨迹};
\end{tikzpicture}
\caption{微团横向匀速移动，纵坐标指数衰减。}\label{fig:c3-h18-q}
\end{figure}
\begin{solution}
随流体微团移动的温度计同时经历场的时间变化与空间输运，固定温度计只经历场在当地的时间变化。先由速度场求图~\ref{fig:c3-h18-q}中的微团轨迹，再代入温度场。


轨迹方程为$\dot x=u,\dot y=-\lambda y$。第一式从初始时刻积分，
\[x-x_0=\int_0^tu\dd s=[us]_0^t=ut.\]
第二式乘积分因子$\ee^{\lambda t}$，乘积法则给出
\[\frac{\dd}{\dd t}\bigl(\ee^{\lambda t}y\bigr)
=\ee^{\lambda t}(\dot y+\lambda y)=0.\]
故$\ee^{\lambda t}y-y_0=\int_0^t0\dd s=0$，即$y=y_0\ee^{-\lambda t}$，其中也包含$y_0=0$的零解。代入温度场，
\[\boxed{T_{\rm parcel}(t)=T_0+ax_0+aut+by_0\ee^{-\lambda t}+ct}.
\]
沿流体的变化率为$\dd T/\dd t=au-\lambda by_0\ee^{-\lambda t}+c$。图~\ref{fig:c3-h18-s}中固定在原位置的温度计只测得$T_{\rm fixed}=T_0+ax_0+by_0+ct$，其变化率为$c$。


当图示参数使初斜率为负、长时斜率为正时，微团温度先降后升。图~\ref{fig:c3-h18-s}中的转折时刻由$\ee^{-\lambda t}=(au+c)/(\lambda by_0)$决定，前提是这个比值在$(0,1)$内。空间输运因而可以改变读数的单调性。

\begin{figure}[H]\centering
\begin{tikzpicture}[>=Stealth,line cap=round,line join=round]
\draw[axis](0,0)--(4.5,0)node[right]{$t$};\draw[axis](0,0)--(0,2.3)node[above]{$T-T_0$};
\draw[curve,domain=0:4,samples=101]plot(\x,{.3*\x+exp(-\x)});
\draw[dashed](0,1)--(4,1.4);
\fill(0,1)circle(1.4pt)node[left]{$1$};
\node[align=left,right]at(4.9,1.2){微团：$0.3t+\ee^{-t}$\\[5pt]固定点：$1+0.1t$};
\end{tikzpicture}
\caption{图示参数在解答中列明。两条曲线从同一读数出发，实线含空间输运产生的额外变化。}\label{fig:c3-h18-s}
\end{figure}

材料导数公式$\partial_tT+\vct v\cdot\nabla T=c+au-\lambda by$与轨迹代入法相同。即使$c=0$，温度场在各固定点不随时间变化，运动微团仍可能感到温度变化；即使局部温度变化非零，沿特定轨迹也可能与空间输运项抵消。


两只温度计在$t=0$的读数相同，之后的差为
\[\boxed{T_{\rm parcel}-T_{\rm fixed}=aut+by_0(\ee^{-\lambda t}-1)}.\]
短时差为$(au-\lambda by_0)t+O(t^2)$，这正是初始空间输运率$\vct v(x_0,y_0)\cdot\nabla T$。解答图采用$x_0=0,by_0=1,\lambda=1,au=0.2,c=0.1$；微团最初经过更低温区域，随后水平输运与显含时间的升温占优，因而读数可以先降后升。


\end{solution}
